% !TeX root = ../Book3.tex
\BookPart{解析函数空间与多复变}{b3:part:07}
\BookPartGuide

前面研究全纯函数时，最常用的收敛方式是在每个紧子集上一致收敛。若还要控制边界大小、比较一个函数与另一个函数的距离，或在给定点值条件下寻找范数最小的函数，就需要把积分结构放进全纯函数的集合中。\chapref{b3:ch:28}选择单位圆盘的两个 Hilbert 空间：\(H^2\) 由圆周上的平方平均控制，\(A^2\) 由圆盘内部的面积积分控制。二者都能用核表示点值，但范数、正交基和边界行为并不相同。

\ProseParagraphBreak

\chapref{b3:ch:29}转向多个复变量。多圆盘上的 Cauchy 公式延续一元的积分方法，Hartogs 现象则展示维数增加后出现的新结论。次调和函数的适当推广、局部 \(\bar\partial\) 方程及复流形语言逐步进入。下图中的联系来自 Bergman 空间：在\chapref{b3:ch:28}熟悉圆盘上的面积范数与再生核以后，可以比较这些构造在多个复变量中怎样变化。
\[
\begin{tikzcd}[row sep=large]
\text{第二十八章：解析函数空间}\arrow[d, "\text{面积范数与核}"]\\
\text{第二十九章：多复变入门中的 Bergman 部分}
\end{tikzcd}
\]
Taylor 系数提供了比较两个空间的具体方法：圆周积分与面积积分赋予同一个单项式不同的权重，点值估计和再生核也随之改变。到了多个复变量，固定其他坐标作一元积分仍然有效，但不同复方向必须彼此相容；这种相容性将产生一元理论中没有的延拓现象。

\ProseParagraphBreak

两章均为选读内容。阅读\chapref{b3:ch:28}时，可以先完成 \(H^2\) 的系数与边界表示，再把 \(A^2\) 的平行计算作比较；零点分解一节解释边界范数为何允许精确除去 Blaschke 因子。继续研究一般 \(H^p\) 空间时，可读 John B. Conway 的《Functions of One Complex Variable II》\cite[pp.~269--299]{Conway1995Functions}中题为“Hardy Spaces on the Disk”的一章。

% !TeX root = ../../Book3.tex
\OptionalChapter{解析函数空间}{b3:ch:28}
\BookChapterNumber{b3:ch:28}
% A26：圆盘 H2、A2 的系数、边界值、零点与再生核；不采用一般 Hp 因子分解。

\begin{chapterabstract}
本章把圆盘上的全纯函数组织为两个可具体计算的 Hilbert 空间。Hardy 空间 \(H^2\) 的范数来自径向圆周平均，其 Taylor 系数平方可和，并决定一个唯一的 \(L^2\) 边界函数；Blaschke 乘积提取全部零点而不改变剩余因子的范数。Bergman 空间 \(A^2\) 则在圆盘内部作面积积分，系数权重随之改变。两种再生核都从正交展开中直接求出，既给出精确的点值估计，也把插值与最佳逼近化成 Hilbert 空间的问题。
\end{chapterabstract}

\begin{learninggoals}
能够区分圆周范数与面积范数，按给定归一化计算单项式的长度及完整系数范数。能够从平方可和系数构造全纯函数和 \(L^2\) 边界值，说明均方收敛、局部一致收敛与边界逐点收敛的差别；能够追踪无限 Blaschke 乘积的内函数性质及除零后范数不变的理由。能够由标准正交基建立 Hardy 核与 Bergman 核，求点值泛函的范数，并用 Bergman 投影求解面积最小二乘问题。
\end{learninggoals}

全纯函数通常由局部信息决定：一个小圆盘内的 Taylor 系数已经确定连通域中的函数。但“确定函数”还没有回答稳定性问题。若只知道函数的某个积分平均不大，它在指定点能有多大？两个全纯函数的系数接近时，它们的边界值是否也接近？给定若干点值，哪一个函数具有最小的积分范数？这些问题需要同时使用复分析与 Hilbert 空间中的正交性。

\ProseParagraphBreak

本章固定单位圆盘 \(\mathbb D=\{z\in\mathbb C:|z|<1\}\) 及单位圆周 \(\mathbb T=\partial\mathbb D\)。圆周测度为 \(\mathrm dm=\mathrm d\theta/(2\pi)\)，面积测度为 \(\mathrm dA=\mathrm dx\,\mathrm dy/\pi\)，二者的总质量都是一。因此常函数一在两个空间中的范数均为一；核公式中是否出现 \(\pi\)，正取决于这项归一化。所有复内积对第一变量线性、对第二变量共轭线性。

\ProseParagraphBreak

先修的复分析内容是\chapref{b3:ch:02}的 Cauchy 公式与 Taylor 展开、\chapref{b3:ch:04}的局部一致全纯极限，以及\chapref{b3:ch:18}的 Jensen 公式和 Blaschke 乘积。积分方面需要控制收敛、Fatou 引理、平方可积空间的完备性和正交展开；用到 Fourier 完备性时，\secref{b3:sec:2802}将明确指出第二卷的结果及角变量转换。这里的 \(H^2\) 是圆盘解析 Hardy 空间，并非 Sobolev 二阶空间，也不是第二卷实变量原子 Hardy 空间的另一种范数写法。

\ProseParagraphBreak

前两节先建立系数、边界和点值之间的对应，第三节研究零点。末两节可以在读完第一节后接读：它们把圆周积分换成面积积分，独立完成完备性和投影，再在最后与 \(H^2\) 比较。我们完整建立的是 \(L^2\) 径向边界值；一般几乎处处非切向边界极限、全部 \(H^p\) 因子分解和移位不变子空间分类均不作为本章证明的前提。若希望从这里继续研究更广的 Hardy 空间，Conway 的《Functions of One Complex Variable II》\cite[pp.~269--299]{Conway1995Functions}的“Hardy Spaces on the Disk”一章讨论函数结构及其与乘法不变子空间的联系；本章先把平方可积情形中的具体构造讲清。

% !TeX root = ../../Book3.tex
\section{\texorpdfstring{\(H^2\)}{H²} 的系数与范数}
\label{b3:sec:2801}
全纯性给出局部 Taylor 展开，却不限制函数接近单位圆周时能增长多快。例如，\(1/(1-z)\) 在单位圆盘内全纯，但在边界点 \(1\) 附近越来越大。若要把全纯函数组成一个适合取极限、作正交投影的空间，便需要一种同时控制所有内部圆周的大小条件。平方平均尤其适合这个目的：幂函数在圆周上正交，积分中的交叉项消失，函数的大小便能够由系数的平方和精确表示。

\subsection{径向平方平均与 Hardy 范数}
\label{b3:subsec:280101}

以下把单位圆盘及其边界分别记为
\[
 \mathbb D\coloneqq\{z\in\mathbb C:|z|<1\},
 \quad \mathbb T\coloneqq\{\zeta\in\mathbb C:|\zeta|=1\}.
\]
圆周采用总质量为一的弧长测度 \(m\)，即当 \(\zeta=e^{\mathrm i\theta}\) 时，\(\mathrm dm(\zeta)=\mathrm d\theta/(2\pi)\)。因此常数函数 \(1\) 的平方积分为一。本章的 \(\mathbb T\) 是复平面中的单位圆；用角变量或模 \(2\pi\) 的实数表示它，积分归一化不变。

\begin{definition}[圆盘上的 Hardy 空间]\label{b3:def:28-hardy-space}
设 \(f:\mathbb D\to\mathbb C\) 全纯。对 \(0<r<1\)，记
\[
 M_2(r,f)^2\coloneqq
 \int_{\mathbb T}|f(r\zeta)|^2\,\mathrm dm(\zeta).
\]
若
\[
 \|f\|_{H^2}^2\coloneqq\sup_{0<r<1}M_2(r,f)^2<\infty,
\]
则称 \(f\) 属于单位圆盘上的\term[hardy-kongjian]{Hardy 空间}[Hardy space]，记为 \(H^2(\mathbb D)\)。这里的 \(H^2\) 专指这一全纯函数空间，上标 \(2\) 表示平方平均，不是二阶 Sobolev 空间。
\end{definition}
\symidx{\(H^2(\mathbb D)\)：圆盘上的 Hardy 空间}
\symidx{\(M_2(r,f)\)：全纯函数的径向平方平均}

定义先对每个内部半径积分，再要求这些积分有一个共同上界。每个固定 \(r<1\) 上，函数连续于紧圆周，积分本来就是有限的；这并不保证 \(r\uparrow1\) 时仍有界。定义也没有预先写 \(f(\zeta)\)，因为一般的圆盘全纯函数尚未在 \(\mathbb T\) 上定义。边界函数必须由内部函数构造，不能先假定它存在。

\ProseParagraphBreak
若全纯函数 \(b\) 在 \(\mathbb D\) 上有界，置
\[
 \|b\|_\infty\coloneqq\sup_{z\in\mathbb D}|b(z)|.
\]
这类函数组成的空间记为 \(H^\infty(\mathbb D)\)。由 \(m(\mathbb T)=1\)，立即得到
\[
 M_2(r,b)\le\|b\|_\infty,
 \quad \|b\|_{H^2}\le\|b\|_\infty.
\]
于是有界全纯函数总属于 \(H^2\)。反方向不成立；下面的系数判据将给出一个在边界附近无界、但仍具有有限平方平均的具体例子。
\symidx{\(H^\infty(\mathbb D)\)：圆盘上的有界全纯函数空间}

\subsection{Taylor 系数与平方可和性}
\label{b3:subsec:280102}

把 \(f\) 在原点展开为 Taylor 级数。先固定一个严格小于一的半径，再作圆周积分，便能避免直接在收敛圆周上代入幂级数的问题。

\begin{proposition}[径向平均的系数公式]\label{b3:prop:28-radial-coefficients}
设 \(f\) 全纯于 \(\mathbb D\)，其 Taylor 展开为 \(f(z)=\sum_{n=0}^{\infty}a_nz^n\)。对每个 \(0<r<1\)，有
\begin{equation}\label{b3:eq:28-radial-coefficients}
 \int_{\mathbb T}|f(r\zeta)|^2\,\mathrm dm(\zeta)
 =\sum_{n=0}^{\infty}|a_n|^2r^{2n}.
\end{equation}
函数 \(r\mapsto M_2(r,f)^2\) 单调不减，且
\begin{equation}\label{b3:eq:28-hardy-coefficient-norm}
 \sup_{0<r<1}M_2(r,f)^2
 =\lim_{r\uparrow1}M_2(r,f)^2
 =\sum_{n=0}^{\infty}|a_n|^2,
\end{equation}
其中允许两端同时为 \(+\infty\)。
\end{proposition}
\begin{proof}
对整数 \(n,k\ge0\)，直接积分得
\[
 \int_{\mathbb T}\zeta^n\overline{\zeta^k}\,\mathrm dm(\zeta)
 =\frac1{2\pi}\int_0^{2\pi}e^{\mathrm i(n-k)\theta}\,\mathrm d\theta
 =\begin{cases}1,&n=k,\\0,&n\ne k.\end{cases}
\]
因而有限多项式 \(p_N(z)=\sum_{n=0}^Na_nz^n\) 满足
\[
 \int_{\mathbb T}|p_N(r\zeta)|^2\,\mathrm dm(\zeta)
 =\sum_{n=0}^N|a_n|^2r^{2n}.
\]
固定 \(r<1\) 时，Taylor 级数在 \(|z|\le r\) 上一致收敛。于是 \(p_N(r\zeta)\) 一致有界，且其模平方一致趋于 \(|f(r\zeta)|^2\)；有限测度上的积分可以取极限。这就证明了\eqref{b3:eq:28-radial-coefficients}。

右边每一项都是 \(r\) 的非负单调函数，因此整个和单调不减。取任意 \(r_j\uparrow1\)，对计数测度使用单调收敛，便有
\[
 \lim_{j\to\infty}\sum_{n=0}^{\infty}|a_n|^2r_j^{2n}
 =\sum_{n=0}^{\infty}|a_n|^2.
\]
单调性保证这个极限也等于所有半径的上确界，故得到第二式。这里没有先假设系数平方可和。
\end{proof}

由此，\(f\in H^2\) 的充要条件恰为 Taylor 系数平方可和。几个可直接计算的例子可以帮助理解这一条件。

\ProseParagraphBreak
对于多项式 \(p(z)=1+2z-z^3\)，有
\[
 M_2(r,p)^2=1+4r^2+r^6,
 \quad \|p\|_{H^2}^2=6.
\]
积分没有保留不同次数之间的交叉项。对于 \(f(z)=1/(1-z)\)，所有系数都为一，所以
\[
 M_2(r,f)^2=\frac1{1-r^2}\longrightarrow\infty.
\]
这说明圆盘内全纯本身远不足以保证 \(H^2\) 条件。

\ProseParagraphBreak
另一方面，函数
\[
 h(z)=\sum_{n=1}^{\infty}\frac{z^n}{n}=-\log(1-z)
\]
属于 \(H^2\)，因为 \(\sum_{n\ge1}n^{-2}<\infty\)。这里的对数分支由在原点取零的幂级数确定；也可注意 \(1-z\) 的实部为正。沿实半径有 \(h(r)=-\log(1-r)\to\infty\)，所以 \(h\notin H^\infty\)。\(H^2\) 控制的是圆周上的总平方平均，容许在逐渐缩小的边界邻域内出现较大值。

\subsection{系数序列的等距表示}
\label{b3:subsec:280103}

系数判据不仅能检验已有函数，还应能反过来构造函数。给定一个平方可和序列，为什么它的幂级数在整个单位圆盘内全纯？这里需要把序列条件转换成紧圆盘上的一致收敛估计。

\begin{theorem}[Hardy 空间的系数表示]\label{b3:thm:28-hardy-coefficients}
设 \(\mathbb N_0=\{0,1,2,\ldots\}\)，并给 \(\ell^2(\mathbb N_0;\mathbb C)\) 配备范数 \(\|(a_n)\|_2^2=\sum_{n\ge0}|a_n|^2\)。映射
\[
 \mathcal C:H^2(\mathbb D)\longrightarrow\ell^2(\mathbb N_0;\mathbb C),
 \quad f\longmapsto\left(\frac{f^{(n)}(0)}{n!}\right)_{n\ge0}
\]
是满射线性等距映射。其逆映射把序列 \((a_n)\) 送到
\[
 f(z)=\sum_{n=0}^{\infty}a_nz^n,
 \quad z\in\mathbb D.
\]
这个级数在每个内部闭圆盘上绝对一致收敛，且所得函数满足 \(\|f\|_{H^2}^2=\sum_{n\ge0}|a_n|^2\)。
\end{theorem}
\begin{proof}
若 \(f\in H^2\)，\cref{b3:prop:28-radial-coefficients}给出等距性。Taylor 系数的唯一性给出单射性，线性则来自逐阶求导的线性。

反之，给定 \((a_n)\in\ell^2\)。对 \(0<R<1\)，序列的 Cauchy--Schwarz 不等式给出
\[
 \sum_{n=0}^{\infty}|a_n|R^n
 \le\left(\sum_{n=0}^{\infty}|a_n|^2\right)^{1/2}
      \left(\sum_{n=0}^{\infty}R^{2n}\right)^{1/2}
 =\frac{\|(a_n)\|_2}{\sqrt{1-R^2}}.
\]
因此幂级数在 \(|z|\le R\) 上绝对一致收敛。由于 \(R<1\) 任意，它定义一个圆盘全纯函数，Taylor 系数就是原序列。再次应用\eqref{b3:eq:28-hardy-coefficient-norm}，便知其 \(H^2\) 范数有限并等于序列范数。满射性得证。
\end{proof}

上述估计还说明：\(H^2\) 范数中的收敛必然蕴含内部紧集上的一致收敛。具体地，若 \(f_j,f\in H^2\) 且 \(\|f_j-f\|_{H^2}\to0\)，则对每个 \(R<1\)，
\begin{equation}\label{b3:eq:28-hardy-compact-control}
 \sup_{|z|\le R}|f_j(z)-f(z)|
 \le\frac{\|f_j-f\|_{H^2}}{\sqrt{1-R^2}}
 \longrightarrow0.
\end{equation}
半径必须固定在一以内；当 \(R\uparrow1\) 时，常数会增大，不能据此推出闭圆盘上一致收敛。

\ProseParagraphBreak
逆命题甚至对多项式列也失败。取 \(f_j(z)=z^j\)，则 \(f_j\to0\) 在每个内部紧集上一致成立，但 \(\|f_j\|_{H^2}=1\)。逐渐靠近边界时仍保留的质量，不能由任何一个固定内部紧集看见。因此局部一致拓扑与 \(H^2\) 范数拓扑需要分别使用。

\subsection{Hilbert 空间结构}
\label{b3:subsec:280104}

系数表示使范数的代数性质与完备性都变得可核算。约定复内积对第一变量线性、对第二变量共轭线性；这一约定将决定后面的核函数在哪个变量上出现共轭。

\begin{proposition}[内积、完备性与多项式逼近]\label{b3:prop:28-hardy-hilbert}
设 \(f,g\in H^2(\mathbb D)\)，其 Taylor 展开分别为 \(f(z)=\sum_{n\ge0}a_nz^n\)、\(g(z)=\sum_{n\ge0}b_nz^n\)。定义
\begin{equation}\label{b3:eq:28-hardy-inner-product}
 \langle f,g\rangle_{H^2}\coloneqq
 \sum_{n=0}^{\infty}a_n\overline{b_n}.
\end{equation}
这个级数绝对收敛，给 \(H^2(\mathbb D)\) 一个复 Hilbert 空间结构，其范数等于径向平均定义的范数。幂函数系 \((z^n)_{n\ge0}\) 是其标准正交基。若 \(p_N(z)=\sum_{n=0}^Na_nz^n\)，则
\[
 \|f-p_N\|_{H^2}^2=\sum_{n>N}|a_n|^2\longrightarrow0.
\]
此外，
\[
 \langle f,g\rangle_{H^2}
 =\lim_{r\uparrow1}
   \int_{\mathbb T}f(r\zeta)\cdot\overline{g(r\zeta)}\,\mathrm dm(\zeta).
\]
\end{proposition}
\begin{proof}
序列的 Cauchy--Schwarz 不等式给出
\[
 \sum_{n=0}^{\infty}|a_n|\cdot|b_n|
 \le\|(a_n)\|_2\cdot\|(b_n)\|_2<\infty.
\]
于是内积良定义，其线性、共轭对称与正定性来自系数运算；若内积平方范数为零，所有系数为零，故全纯函数恒为零。

为看清完备性如何保留全纯性，设 \((f_j)\) 是 \(H^2\) Cauchy 列，系数列记为 \(a^{(j)}\)。等距性使 \((a^{(j)})\) 成为 \(\ell^2\) Cauchy 列，因而收敛到某个 \(a\in\ell^2\)。序列空间完备性也可直接检验：各坐标分别有极限，有限平方和取极限给 \(\sum_{n=0}^N|a_n|^2\le\sup_j\|a^{(j)}\|_2^2\)，再令 \(N\to\infty\) 得到 \(a\in\ell^2\)；对 Cauchy 误差作同样的有限和取极限，得到 \(\|a^{(j)}-a\|_2\to0\)。

由\cref{b3:thm:28-hardy-coefficients}，\(a\) 构造出一个全纯函数 \(f\in H^2\)，并且 \(\|f_j-f\|_{H^2}=\|a^{(j)}-a\|_2\to0\)。所以极限确实仍在原函数空间内。幂函数的正交性、单位范数及部分和余项公式，也直接来自系数表示。

最后，先对有限多项式使用圆周正交性，再由内部圆周上的一致收敛取极限，得到
\[
 \int_{\mathbb T}f(r\zeta)\cdot\overline{g(r\zeta)}\,\mathrm dm(\zeta)
 =\sum_{n=0}^{\infty}a_n\overline{b_n}r^{2n}.
\]
右边各项由可和数列 \(|a_n|\cdot|b_n|\) 控制。对计数测度使用支配收敛，就得到所述内积极限。
\end{proof}

这里的标准正交基由解析多项式组成，只含非负幂。它不同于全部圆周 \(L^2\) 空间的三角基：后者还含负频率。下一节将把这一区别解释为边界数据能够来自圆盘全纯函数的精确条件。

% !TeX root = ../../Book3.tex
\section{\texorpdfstring{\(L^2\)}{L²} 边界值与 Hardy 核}
\label{b3:sec:2802}
径向平方平均的有界性已经等价于系数平方可和。现在把系数放回圆周：每个内部圆周上的 \(f(r\zeta)\) 都含有衰减因子 \(r^n\)，让 \(r\uparrow1\) 就是在逐渐撤去这种衰减。平方可和性保证这个过程在 \(L^2\) 范数中有极限。所得边界函数将把内部点值表示为一个连续线性泛函，核函数由此出现。

\subsection{径向极限与边界函数}
\label{b3:subsec:280201}

对 \(f\in H^2(\mathbb D)\)，把圆周函数 \(\zeta\mapsto f(r\zeta)\) 简记为 \(f_r\)。不同的 \(r\) 对应不同的内部圆周，但参数都放在同一个测度空间 \((\mathbb T,m)\) 上，这样才可以谈论它们的 \(L^2\) 距离。

\begin{theorem}[Hardy 函数的平方积分边界值]\label{b3:thm:28-hardy-boundary}
设 \(f\in H^2(\mathbb D)\)，其 Taylor 系数为 \((a_n)_{n\ge0}\)，圆周测度为 \(\mathrm dm(e^{\mathrm i\theta})=\mathrm d\theta/(2\pi)\)。存在唯一的 \(f^*\in L^2(\mathbb T,m;\mathbb C)\)，使
\begin{equation}\label{b3:eq:28-hardy-radial-limit}
 \lim_{r\uparrow1}\|f_r-f^*\|_{L^2(m)}=0.
\end{equation}
这里的唯一性按几乎处处相等理解。具体地，
\[
 f^*=\lim_{N\to\infty}\sum_{n=0}^Na_n\zeta^n
 \quad\text{于 }L^2(\mathbb T,m),
\]
并且
\begin{equation}\label{b3:eq:28-boundary-error}
 \|f_r-f^*\|_{L^2(m)}^2
 =\sum_{n=0}^{\infty}|a_n|^2(1-r^n)^2,
 \quad \|f^*\|_{L^2(m)}=\|f\|_{H^2}.
\end{equation}
对任意 \(f,g\in H^2(\mathbb D)\)，还有
\[
 \langle f,g\rangle_{H^2}
 =\int_{\mathbb T}f^*(\zeta)\cdot\overline{g^*(\zeta)}\,\mathrm dm(\zeta).
\]
称 \(f^*\) 为 \(f\) 的 \(L^2\) \term[jingxiang-bianjiezhi]{径向边界值}[radial boundary value]。
\end{theorem}
\symidx{\(f^*\)：Hardy 函数的 \(L^2\) 径向边界值}
\begin{proof}
令 \(F_N(\zeta)=\sum_{n=0}^Na_n\zeta^n\)。圆周正交性给出
\[
 \|F_N-F_M\|_{L^2(m)}^2=\sum_{n=M+1}^N|a_n|^2
 \quad(N>M).
\]
因为 \((a_n)\in\ell^2\)，\((F_N)\) 是 \(L^2\) Cauchy 列。平方可积函数空间的完备性给出极限 \(f^*\)。令 \(N\to\infty\)，有限和的平方范数公式同时给出 \(\|f^*\|_2^2=\sum|a_n|^2\)。这一构造只用非负幂的正交性，尚不需要全部三角系的完备性。

固定 \(r<1\)。\(f_r\) 是 \(\sum_{n\ge0}a_nr^n\zeta^n\) 在圆周上的一致极限，因而也是 \(L^2\) 极限。两列有限和相减并取范数极限，得到\eqref{b3:eq:28-boundary-error}的第一式。对每个 \(n\)，\((1-r^n)^2\to0\)，且它不超过一；可和控制来自 \(\sum|a_n|^2<\infty\)。对计数测度使用支配收敛，便得\eqref{b3:eq:28-hardy-radial-limit}。

范数空间中的极限唯一，所以 \(f^*\) 作为 \(L^2\) 元素唯一。对两个有限和计算内积，再以 Cauchy--Schwarz 控制取极限的误差，得到
\[
 \int_{\mathbb T}f^*\overline{g^*}\,\mathrm dm
 =\sum_{n=0}^{\infty}a_n\overline{b_n}
 =\langle f,g\rangle_{H^2}.
\]
这也表明边界映射是线性等距映射。
\end{proof}

若 \(f\) 连续到闭圆盘，那么 \(f_r\) 一致趋于通常的边界限制 \(f|_{\mathbb T}\)，所以上述 \(f^*\) 就是它的 \(L^2\) 等价类。一般情况下，星号表示通过极限得到的数据，不表示先在圆周上任意给 \(f\) 补值。

\ProseParagraphBreak
需要区分范数极限与逐点径向极限。\eqref{b3:eq:28-hardy-radial-limit}本身只直接保证存在几乎处处收敛的半径子列：选 \(r_j\uparrow1\)，使 \(\|f_{r_j}-f^*\|_2^2\le2^{-j}\)，则 Tonelli 定理给出
\[
 \int_{\mathbb T}\sum_{j=1}^{\infty}|f(r_j\zeta)-f^*(\zeta)|^2
 \,\mathrm dm(\zeta)\le\sum_{j=1}^{\infty}2^{-j}<\infty.
\]
于是误差平方和在几乎处处有限，故 \(f(r_j\zeta)\to f^*(\zeta)\) 几乎处处成立。这里得到沿所选半径子列的几乎处处收敛；所有 \(r\uparrow1\) 的几乎处处径向极限需要进一步的边界定理。

\ProseParagraphBreak
特别地，把有界全纯因子乘进来时，必须说明边界乘积也确实由内部乘积得到。两个 \(L^2\) 极限通常只能直接相乘到 \(L^1\)；这里额外的有界性同时控制了乘积范数与边界值。

\begin{proposition}[有界全纯因子的乘法]\label{b3:prop:28-hardy-multiplier}
设 \(b\in H^\infty(\mathbb D)\)、\(f\in H^2(\mathbb D)\)，并取归一化圆周测度 \(m\)。则 \(bf\in H^2(\mathbb D)\)，且
\[
 \|bf\|_{H^2}\le\|b\|_\infty\cdot\|f\|_{H^2},
 \quad (bf)^*=b^*f^*\quad m\text{-几乎处处}.
\]
此外，\(|b^*|\le\|b\|_\infty\) 几乎处处。若还满足 \(|b^*|=1\) 几乎处处，则
\[
 \|bf\|_{H^2}=\|f\|_{H^2}.
\]
\end{proposition}
\begin{proof}
对每个 \(0<r<1\)，内部的点态界给出
\[
 \int_{\mathbb T}|b(r\zeta)\cdot f(r\zeta)|^2\,\mathrm dm(\zeta)
 \le\|b\|_\infty^2
       \int_{\mathbb T}|f(r\zeta)|^2\,\mathrm dm(\zeta).
\]
取上确界即得 \(bf\in H^2\) 及范数界，且 \(b\in H^2\) 已由有界性保证。对 \(b,f,bf\) 的三份 \(L^2\) 径向误差同时选取半径 \(r_j\uparrow1\)，使它们的平方范数之和小于 \(2^{-j}\)。刚才的可和误差论证说明，除去一个零测集后，三个函数沿这同一子列都收敛到各自的边界函数。在这些圆周点上取通常复数乘法的极限，便有 \((bf)^*=b^*f^*\)。同样，\(|b(r_j\zeta)|\le\|b\|_\infty\) 给出边界模长界。

若 \(|b^*|=1\) 几乎处处，边界等距性给出
\[
 \|bf\|_{H^2}^2
 =\int_{\mathbb T}|b^*|^2\cdot|f^*|^2\,\mathrm dm
 =\int_{\mathbb T}|f^*|^2\,\mathrm dm
 =\|f\|_{H^2}^2.
\]
这一子列上的几乎处处收敛使极限乘积与边界函数相对应。
\end{proof}

例如乘以 \(z\) 只把系数向后移动一位，因而是 \(H^2\) 上的等距映射；边界解释则是 \(|\zeta|=1\)。这个映射不是满射，因为所得函数都在原点为零。范数不变并不意味着不丢失可取到的函数。

\subsection{边界 Fourier 系数}
\label{b3:subsec:280202}

在 \(L^2(\mathbb T,m)\) 中，每个函数也属于 \(L^1\)，因为 \(m(\mathbb T)=1\)。因此对任意整数 \(k\)，其 Fourier 系数
\[
 \widehat F(k)\coloneqq
 \int_{\mathbb T}F(\zeta)\cdot\overline{\zeta^k}\,\mathrm dm(\zeta)
\]
都定义良好。\(\widehat F(k)\) 关于 \(F\) 是连续线性泛函，且 \(|\widehat F(k)|\le\|F\|_2\)。前面的构造表明，解析函数的边界系数只能出现在 \(k\ge0\) 的位置。

\ProseParagraphBreak
这里复用第二卷定理3.3.2与推论3.3.3：在总质量为一的圆周上，\((e^{2\pi\mathrm i kt})_{k\in\mathbb Z}\) 是 \(L^2\) 的标准正交基，Fourier 系数映射是到 \(\ell^2(\mathbb Z)\) 的满射等距映射。令 \(\theta=2\pi t\) 后，基函数变为 \(\zeta^k\)，测度变为 \(\mathrm d\theta/(2\pi)\)，故 Parseval 等式的常数仍为一。这些结论指 \(L^2\) 展开，不要求 Fourier 级数逐点收敛。

\begin{proposition}[边界函数的频率刻画]\label{b3:prop:28-hardy-boundary-range}
取单位圆周 \(\mathbb T\) 及测度 \(\mathrm dm(e^{\mathrm i\theta})=\mathrm d\theta/(2\pi)\)，定义闭子空间
\[
 L^2_+(\mathbb T)\coloneqq
 \{F\in L^2(\mathbb T,m;\mathbb C):\widehat F(k)=0\text{ 对所有 }k<0\}.
\]
边界映射
\[
 H^2(\mathbb D)\longrightarrow L^2_+(\mathbb T),
 \quad f\longmapsto f^*
\]
是满射线性等距映射。若 \(f(z)=\sum_{n\ge0}a_nz^n\)，则
\[
 \widehat{f^*}(k)=
 \begin{cases}a_k,&k\ge0,\\0,&k<0.\end{cases}
\]
任意 \(F\in L^2(\mathbb T,m;\mathbb C)\) 到该子空间的正交投影为
\[
 P_+F=\sum_{n=0}^{\infty}\widehat F(n)\zeta^n
 \quad\text{于 }L^2(\mathbb T,m).
\]
特别地，\(\|P_+F\|_2\le\|F\|_2\)。
\end{proposition}
\symidx{\(L^2_+(\mathbb T)\)：负 Fourier 系数为零的圆周闭子空间}
\symidx{\(P_+\)：圆周 \(L^2\) 到非负频率子空间的正交投影}
\begin{proof}
每个条件 \(\widehat F(k)=0\) 都是一个连续线性泛函的核，因此这些核的交是闭线性子空间。对 \(f^*\) 的多项式逼近取 Fourier 系数，系数泛函的连续性给出所列公式，故边界映射的像包含于 \(L^2_+\)。

反之，设 \(F\in L^2_+\)。由第二卷的 Parseval 等式，\((\widehat F(n))_{n\ge0}\) 平方可和，且
\[
 \sum_{n=0}^{\infty}|\widehat F(n)|^2=\|F\|_2^2.
\]
系数表示定理构造出 \(f(z)=\sum_{n\ge0}\widehat F(n)z^n\in H^2\)。\(f^*\) 与 \(F\) 的所有整数频率系数相同；三角系的完备性说明二者在 \(L^2\) 中相等。于是满射性成立。

对一般 \(F\)，Parseval 等式也保证只保留非负频率的级数在 \(L^2\) 中收敛，所得函数属于 \(L^2_+\)。余项 \(F-P_+F\) 只含负频率，因而与每个非负幂正交，再由非负幂在 \(L^2_+\) 中的稠密性知，它与整个子空间正交。这就是正交投影，且
\[
 \|F\|_2^2=\|P_+F\|_2^2+\|F-P_+F\|_2^2
\]
给出其范数界。
\end{proof}

例如 \(F(\zeta)=\overline\zeta\) 的全部质量在频率 \(-1\)，故 \(P_+F=0\)。它不能是任何圆盘 \(H^2\) 函数的边界值，尽管它光滑且平方可积。又如 \(F(\zeta)=\operatorname{Re}\zeta=(\zeta+\overline\zeta)/2\)，其正交投影为 \(\zeta/2\)，不是原数据本身。边界数据的可积性与它能否由全纯函数产生，是两个不同条件。

\ProseParagraphBreak
还有一个有用的刚性现象：若 \(f\in H^2\) 的边界值实值几乎处处，则 \(f\) 必为实常数。因为实值 \(F\) 满足 \(\widehat F(-n)=\overline{\widehat F(n)}\)，负频率消失就迫使全部非零频率消失。这里结论依赖边界函数来自 \(H^2\)；仅给某个全纯函数零测集以外的边界行为，不能不加增长条件就代替这项假设。

\subsection{Cauchy 表示与点值控制}
\label{b3:subsec:280203}

经典 Cauchy 公式在内部圆周上成立，那里边值连续。本节的边界函数只有平方可积性。要把圆周移到 \(\mathbb T\)，应先检验积分核的可积性与极限，而不是把经典公式中的半径直接换成一。

\begin{proposition}[边界 Cauchy 表示与最优点值界]\label{b3:prop:28-hardy-cauchy}
设 \(f\in H^2(\mathbb D)\)，\(f^*\) 是其相对于归一化圆周测度 \(m\) 的 \(L^2\) 边界值。对任意 \(w\in\mathbb D\)，有
\begin{equation}\label{b3:eq:28-hardy-cauchy}
 f(w)=\int_{\mathbb T}
       \frac{f^*(\zeta)}{1-w\overline\zeta}\,\mathrm dm(\zeta),
\end{equation}
且
\begin{equation}\label{b3:eq:28-hardy-point-evaluation}
 |f(w)|\le\frac{\|f\|_{H^2}}{\sqrt{1-|w|^2}}.
\end{equation}
对每个固定 \(w\)，右端常数是最优的；除零函数外，等号成立当且仅当
\[
 f(z)=\frac{c}{1-\overline w z}
 \quad(z\in\mathbb D),\quad c\in\mathbb C\setminus\{0\}.
\]
\end{proposition}
\begin{proof}
固定 \(w\in\mathbb D\)。由于 \(|1-w\overline\zeta|\ge1-|w|>0\)，积分核有界，而 \(f^*\in L^1(m)\)，所以积分绝对存在。几何级数
\[
 \frac1{1-w\overline\zeta}
 =\sum_{n=0}^{\infty}w^n\overline\zeta^{\,n}
\]
在 \(\zeta\in\mathbb T\) 上绝对一致收敛。乘以 \(f^*\) 后，用可积控制 \(|f^*|/(1-|w|)\) 交换积分与求和，即得
\[
 \int_{\mathbb T}\frac{f^*(\zeta)}{1-w\overline\zeta}\,\mathrm dm
 =\sum_{n=0}^{\infty}w^n\widehat{f^*}(n)
 =\sum_{n=0}^{\infty}a_nw^n=f(w).
\]
这个证明同时解释了为什么表示式只读取非负 Fourier 系数。

函数 \(K_w(\zeta)=1/(1-\overline w\zeta)\) 的圆周平方范数由几何级数和正交性给出：
\[
 \|K_w\|_{L^2(m)}^2
 =\sum_{n=0}^{\infty}|w|^{2n}
 =\frac1{1-|w|^2}.
\]
于是\eqref{b3:eq:28-hardy-cauchy}是 \(\langle f^*,K_w\rangle_{L^2}\)，Cauchy--Schwarz 不等式与边界等距性给出点值界。对非零 \(f\)，等号成立当且仅当 \(f^*=cK_w\)；边界映射的单射性使这等价于 \(f(z)=c/(1-\overline w z)\)。这些函数确实属于 \(H^2\)，所以常数不能减小。
\end{proof}

若 \(f\) 连续到闭圆盘，令 \(\zeta=e^{\mathrm i\theta}\) 并利用 \(\mathrm d\zeta=\mathrm i\zeta\,\mathrm d\theta\)，就会把\eqref{b3:eq:28-hardy-cauchy}写成通常的 \((2\pi\mathrm i)^{-1}\int_{\mathbb T}f(\zeta)/(\zeta-w)\,\mathrm d\zeta\)。对于一般 \(f\in H^2\)，Lebesgue 积分中的边界函数是一个 \(L^2\) 等价类，改变零测集上的值不会改变积分。

\ProseParagraphBreak
同一个边界函数还有 Poisson 表示。以本卷\secref{b3:sec:1601}的单位圆盘归一化记核
\[
 P(w,\zeta)=\frac{1-|w|^2}{|\zeta-w|^2}
 =1+\sum_{n\ge1}\bigl(w^n\overline\zeta^{\,n}
                     +\overline w^{\,n}\zeta^n\bigr).
\]
对固定 \(w\in\mathbb D\)，该级数在圆周上一致绝对收敛，故它也可与 \(f^*\in L^1\) 逐项积分。含 \(\overline w^{\,n}\) 的部分乘以 \(\widehat{f^*}(-n)=0\)，从而
\begin{equation}\label{b3:eq:28-hardy-poisson}
 f(w)=\int_{\mathbb T}P(w,\zeta)\cdot f^*(\zeta)\,\mathrm dm(\zeta).
\end{equation}
这只是把本处 \(L^2\) 数据代入一个有界内部核，并没有引用前章只针对连续数据的边界收敛结论。

\ProseParagraphBreak
对于任意 \(F\in L^2\)，Cauchy 积分读取 \(P_+F\)，Poisson 积分则保留正、负两组频率。以 \(F(\zeta)=\overline\zeta\) 为例，前者为零，后者为 \(\overline w\)。只有当负频率消失时，两种内部表示才相同。因此 Poisson 延拓通常是调和的，全纯性还需要边界数据满足频率条件。

\subsection{Hardy 核与再生性质}
\label{b3:subsec:280204}

点值估计说明，在 \(H^2\) 中取某个内点的函数值是一个连续线性操作。Cauchy 表示里的核不仅属于边界的 \(L^2\) 空间，它本身也是一个圆盘全纯函数。由此可以把点值直接表示成两个解析函数的内积。

\begin{definition}[再生核]\label{b3:def:28-reproducing-kernel}
设 \(X\) 是非空集合，\(\mathcal H\) 是由 \(X\) 上复值函数组成的复 Hilbert 空间，内积对第一变量线性。若函数 \(K:X\times X\to\mathbb C\) 满足：对每个 \(w\in X\)，函数 \(K_w=K(\cdot,w)\) 属于 \(\mathcal H\)，且
\[
 f(w)=\langle f,K_w\rangle_{\mathcal H}
 \quad(f\in\mathcal H),
\]
则称 \(K\) 为 \(\mathcal H\) 的\term[zaisheng-he]{再生核}[reproducing kernel]，并称 \(\mathcal H\) 为\term[zaisheng-he-hilbert-kongjian]{再生核 Hilbert 空间}[reproducing kernel Hilbert space]。
\end{definition}

定义中的元素是实际的函数，因而每个点值已经有意义。一般 \(L^2\) 空间的元素只按几乎处处相等区分，单点值不能由其等价类确定；不能把任意 \(L^2\) 空间直接当作再生核空间。本章的 \(H^2\) 元素始终是在盘内逐点定义的全纯函数，边界映射才把它们识别为一个特殊的 \(L^2\) 子空间。

\begin{theorem}[Hardy 核]\label{b3:thm:28-hardy-kernel}
在单位圆盘 \(\mathbb D\) 上，给 \(H^2(\mathbb D)\) 配备内积
\(\langle f,g\rangle=\int_{\mathbb T}f^*\overline{g^*}\,\mathrm dm\)，其中 \(\mathrm dm(e^{\mathrm i\theta})=\mathrm d\theta/(2\pi)\)。其再生核为
\begin{equation}\label{b3:eq:28-hardy-kernel}
 S(z,w)\coloneqq\frac1{1-z\overline w},
 \quad z,w\in\mathbb D.
\end{equation}
对每个 \(w\in\mathbb D\)，记 \(k_w=S(\cdot,w)\)。则对任意 \(f\in H^2(\mathbb D)\)，有
\[
 f(w)=\langle f,k_w\rangle_{H^2},
 \quad \|k_w\|_{H^2}^2=S(w,w)=\frac1{1-|w|^2}.
\]
该核称为圆盘的\term[hardy-he]{Hardy 核}[Hardy kernel]，也称\term[szego-he]{Szegő 核}[Szegő kernel]。
\end{theorem}
\symidx{\(S(z,w)\)：单位圆盘的 Hardy 核}
\begin{proof}
固定 \(w\in\mathbb D\)。\(k_w(z)=\sum_{n\ge0}\overline w^{\,n}z^n\) 的系数平方和为 \((1-|w|^2)^{-1}\)，所以它属于 \(H^2\)。若 \(f(z)=\sum a_nz^n\)，则
\[
 \langle f,k_w\rangle_{H^2}
 =\sum_{n=0}^{\infty}a_n\overline{\overline w^{\,n}}
 =\sum_{n=0}^{\infty}a_nw^n=f(w).
\]
再取 \(f=k_w\) 就得到 \(\|k_w\|^2=k_w(w)=S(w,w)\)。

再生性还保证唯一性：若 \(\widetilde k_w\) 也代表同一点值，则 \(\langle f,k_w-\widetilde k_w\rangle=0\) 对所有 \(f\in H^2\) 成立。取 \(f=k_w-\widetilde k_w\)，可得两函数的距离为零，故相等。
\end{proof}

核在 \(z\) 上全纯，在 \(w\) 上反全纯；\(\overline w\) 的位置由内积的线性约定决定。交换两个变量有 \(S(z,w)=\overline{S(w,z)}\)，而
\[
 \langle k_w,k_z\rangle_{H^2}=k_w(z)=S(z,w).
\]
这些等式也说明，核的两个变量并非两个独立复制的积分变量：第一个变量表示函数的自变量，第二个变量指定要恢复的点值。

\ProseParagraphBreak
核函数还给出一个可直接求解的极值问题。给定 \(w\in\mathbb D\) 与 \(c\in\mathbb C\)，在满足 \(f(w)=c\) 的全部 \(f\in H^2\) 中，唯一范数最小的函数是
\[
 f_0(z)=c(1-|w|^2)\cdot k_w(z),
 \quad \|f_0\|_{H^2}=|c|\sqrt{1-|w|^2}.
\]
确实，\(f_0(w)=c\)；对任一同样满足该约束的 \(f\)，有 \(\langle f-f_0,k_w\rangle=(f-f_0)(w)=0\)，因此
\[
 \|f\|_{H^2}^2=\|f_0\|_{H^2}^2+\|f-f_0\|_{H^2}^2.
\]
最小性与唯一性同时得出。再生核把一个点值约束变成 Hilbert 空间中的正交分解；\secref{b3:sec:2805}会用面积范数得到另一种核，并比较同一点值的代价如何改变。

% !TeX root = ../../Book3.tex
\section{Blaschke 乘积与零点}
\label{b3:sec:2803}

上一节把 \(H^2(\mathbb D)\) 中的解析函数与圆周上的平方可积边界函数联系起来。现在要问：内部的零点怎样影响这份边界资料？直接除以 \(z-a\) 虽能消去一个零点，却会改变圆周各处的模。Blaschke 因子在单位圆周模为一，因而适合把零点从函数中分离出来，同时保留边界模和 Hardy 范数。

\ProseParagraphBreak

\chapref{b3:ch:18}已经证明有界解析函数的零点满足 Blaschke 条件，并从满足该条件的零点列构造了乘积。本节不再重做那套构造，而要补足两件事：平方平均有界为何也给出同一零点条件；无穷乘积为何仍有模为一的边界值，从而可以在 \(H^2\) 中作等范数除法。仍取 \(\mathbb D=\{\,z\in\mathbb C\mid |z|<1\,\}\)，圆周测度为 \(\mathrm dm=\mathrm d\theta/(2\pi)\)，\(f^*\) 表示\cref{b3:thm:28-hardy-boundary}建立的 \(L^2\) 边界函数。

\subsection{零点列与 Blaschke 条件}
\label{b3:subsec:280301}

解析性只禁止零点在圆盘内部聚集；它并未限制零点沿边界聚集得多快。\(H^2\) 的额外约束来自所有内部圆周平方平均的共同上界。Jensen 公式恰好把这个平均界转换为零点的加权总量。

\Needspace{6\baselineskip}
\begin{theorem}[Hardy 函数的零点条件]\label{b3:thm:28-hardy-zero-condition}
设 \(f\in H^2(\mathbb D)\)、\(f\not\equiv0\)，原点是它的 \(m\) 重零点，其中 \(m\ge0\)，并记第一个非零 Taylor 系数为 \(c_m=f^{(m)}(0)/m!\)。把全部非零零点按重数列为 \((a_j)\)，允许有限列或空列。则
\begin{equation}\label{b3:eq:28-zero-budget}
 \sum_j-\log|a_j|
 \le \log\frac{\|f\|_{H^2}}{|c_m|}<\infty,
 \quad
 \sum_j(1-|a_j|)<\infty.
\end{equation}
\end{theorem}
\begin{proof}
固定 \(0<r<1\)。函数 \(f\) 在闭盘 \(|z|\le r\) 的开邻域全纯，所以\cref{b3:thm:18-jensen}适用；即使观察圆周上有零点，其对数模仍可积。该公式给出
\[
 \int\log|f(re^{\mathrm i\theta})|\,\mathrm dm
 =\log|c_m|+m\log r
       +\sum_{|a_j|<r}\log\frac r{|a_j|}.
\]
对数的凹性以及 \(m(\partial\mathbb D)=1\) 给出
\[
 \int\log|f(re^{\mathrm i\theta})|\,\mathrm dm
 \le\frac12\log\int|f(re^{\mathrm i\theta})|^2\,\mathrm dm
 \le\log\|f\|_{H^2}.
\]
这是\eqref{b3:eq:18-mean-power-bound}在指数二时的版本。故
\[
 \sum_j\left(\log\frac r{|a_j|}\right)_+
 \le\log\frac{\|f\|_{H^2}}{|c_m|}-m\log r.
\]
左侧每一项随 \(r\uparrow1\) 非减，趋于 \(-\log|a_j|\)。对计数测度使用单调收敛，且注意 \(m\log r\to0\)，便得到第一个不等式。最后利用 \(1-t\le-\log t\)（\(0<t<1\)）得到 Blaschke 条件。
\end{proof}

原点零点由有限整数 \(m\) 单独记录，不能把 \(-\log0\) 塞入和式。式\eqref{b3:eq:28-zero-budget}还保留了定量信息：若 \(f(0)\ne0\)，则
\[
 \prod_j|a_j|\ge\frac{|f(0)|}{\|f\|_{H^2}}>0.
\]
中心值相对范数越大，零点可以消耗的对数总量就越少。例如 \(|f(0)|=\|f\|_{H^2}\) 时，系数等距已经迫使所有正次系数为零，函数是非零常数；零点预算为零与这个结论一致。

\ProseParagraphBreak

求和必须计入重数。一个点即使非常靠近边界，若把它重复足够多次，仍可使总和发散。这里控制的并非圆盘中的几何点集，而是解析函数所携带的计重数零点列。

\subsection{Blaschke 乘积的收敛}
\label{b3:subsec:280302}

先恢复乘积的归一化。对 \(a\in\mathbb D\setminus\{0\}\)，取
\begin{equation}\label{b3:eq:28-blaschke-factor}
 b_a(z)=\frac{|a|}{a}\frac{a-z}{1-\overline a z},
 \quad b_a(0)=|a|.
\end{equation}
它与\cref{b3:def:06-finite-blaschke}使用的 \((z-a)/(1-\overline a z)\) 只差一个模一常数，所以零点与模都相同。有限乘积中这个常数可以并入整体相位；无穷乘积却必须固定相位，使靠近边界的因子在内部紧集上趋于一。

\ProseParagraphBreak

具体说，设 \((a_j)\subset\mathbb D\setminus\{0\}\) 按有限重数重复，满足 \(\sum_j(1-|a_j|)<\infty\)，另取整数 \(m\ge0\) 和 \(|\eta|=1\)。这个条件本身已保证任何内部闭盘中只有有限项，否则每项缺陷都有共同正下界。\cref{b3:thm:18-blaschke-zeros}因而适用，给出
\begin{equation}\label{b3:eq:28-blaschke-product}
 B(z)=\eta z^m\prod_j b_{a_j}(z).
\end{equation}
乘积在圆盘内局部一致收敛，不恒为零，其零点恰为所列零点，且 \(|B(z)|\le1\)。空乘积取一。回忆收敛的关键估计：若 \(a=\rho e^{\mathrm i\varphi}\)，则
\[
 \begin{aligned}
 1-b_a(z)&=(1-\rho)\frac{1+e^{-\mathrm i\varphi}z}
                              {1-\rho e^{-\mathrm i\varphi}z},\\
 \sup_{|z|\le s}|1-b_a(z)|&\le\frac{1+s}{1-s}(1-|a|)
 \quad(0<s<1).
 \end{aligned}
\]
由此得到尾因子的局部对数级数绝对一致收敛；它的指数无零点，有限头部则保留指定的重数。这个论证也说明不能去掉式\eqref{b3:eq:28-blaschke-factor}中的相位，只凭每个因子模不超过一来断言乘积收敛。

\ProseParagraphBreak

有限乘积在每个边界点模为一，但局部一致极限只控制圆盘内部，不能直接传递这项边界性质。需要结合刚刚建立的 \(L^2\) 边界理论。为此先给出所用的术语。

\Needspace{6\baselineskip}
\begin{definition}\label{b3:def:28-inner}
设 \(I\) 在 \(\mathbb D\) 全纯且有界。若它的 \(L^2(\partial\mathbb D,m)\) 边界函数满足 \(|I^*|=1\) 几乎处处，则称 \(I\) 为\term[nei-hanshu]{内函数}[inner function]。这里的边界函数由 \(I\in H^2(\mathbb D)\) 及径向 \(L^2\) 收敛确定，不要求连续到闭圆盘。
\end{definition}

\Needspace{6\baselineskip}
\begin{theorem}[Blaschke 乘积的内函数性质]\label{b3:thm:28-blaschke-inner}
设 \(m\ge0\) 为整数，\(|\eta|=1\)，且非零点列 \((a_j)\subset\mathbb D\) 按有限重数重复并满足 Blaschke 条件。式\eqref{b3:eq:28-blaschke-product}定义的有限或无穷乘积 \(B\) 是内函数，且
\[
 \|B\|_{H^2}=1,
 \quad
 \lim_{r\uparrow1}\int-\log|B(re^{\mathrm i\theta})|\,\mathrm dm=0.
\]
\end{theorem}
\begin{proof}
由构造，\(B\) 在盘内模不超过一，故属于 \(H^2\)，并存在 \(L^2\) 边界函数 \(B^*\)。Blaschke 条件还给出 \(\sum_j-\log|a_j|<\infty\)：除了有限项以外 \(|a_j|\ge1/2\)，而在 \(1/2\le t<1\) 时有 \(-\log t\le2(1-t)\)。因此 \(B\) 的第一个非零 Taylor 系数为
\[
 c_m=\eta\prod_j|a_j|\ne0.
\]
对每个 \(0<r<1\) 使用 Jensen 公式，得到
\[
 \int\log|B(re^{\mathrm i\theta})|\,\mathrm dm
 =m\log r+\sum_j\log|a_j|
             +\sum_{|a_j|<r}\log\frac r{|a_j|}.
\]
最后一个非负和随 \(r\uparrow1\) 单调趋于 \(\sum_j-\log|a_j|\)。两个有限级数恰好抵消，\(m\log r\to0\)，所以对数平均趋于零。

由径向 \(L^2\) 收敛，可以选 \(r_k\uparrow1\)，使 \(B(r_ke^{\mathrm i\theta})\to B^*(e^{\mathrm i\theta})\) 几乎处处。因内部模不超过一，\(|B^*|\le1\) 几乎处处。对非负函数 \(-\log|B(r_ke^{\mathrm i\theta})|\) 使用 Fatou 引理；在零值处规定 \(-\log0=+\infty\)，得到
\[
 0\le\int-\log|B^*|\,\mathrm dm
 \le\liminf_{k\to\infty}\int-\log|B(r_ke^{\mathrm i\theta})|\,\mathrm dm=0.
\]
故 \(-\log|B^*|=0\) 几乎处处，即 \(|B^*|=1\)。边界等距于是给 \(\|B\|_{H^2}^2=\int|B^*|^2\,\mathrm dm=1\)。
\end{proof}

这里没有把一个无控制的对数极限移进积分。我们先由 Jensen 公式算出非负积分趋于零，再用 Fatou 引理判断极限。所需的逐点收敛也只有 \(L^2\) 收敛所提供的一条子列；无需预先假设所有径向极限都存在。

\ProseParagraphBreak

乘积不仅在内部紧集收敛，在 Hardy 范数下也收敛。设零点列无限，取
\[
 B_N=\eta z^m\prod_{j=1}^N b_{a_j},
 \quad T_N=\prod_{j>N}b_{a_j},
 \quad B=B_NT_N.
\]
两个乘积的边界模均为一。利用\cref{b3:prop:28-hardy-multiplier}的边界乘法公式，以及内积对第一变量线性的约定，有
\[
 \langle B,B_N\rangle_{H^2}
 =\int B_N^*T_N^*\overline{B_N^*}\,\mathrm dm
 =\int T_N^*\,\mathrm dm=T_N(0)=\prod_{j>N}|a_j|.
\]
最后一步是边界 Fourier 零次系数等于解析函数的常数项。因此
\begin{equation}\label{b3:eq:28-blaschke-hardy-convergence}
 \|B-B_N\|_{H^2}^2
 =2\left(1-\prod_{j>N}|a_j|\right)\longrightarrow0.
\end{equation}
尾部对数和趋于零，正是尾乘积在原点趋于一的原因。这项平均意义的逼近不声称在整个闭圆盘上一致收敛。

\subsection{零点分解与 Hardy 范数}
\label{b3:subsec:280303}

对有界函数，第六章和第十八章已用最大模原理证明除以 Blaschke 因子仍有相同模界。\(H^2\) 函数可能无界，此时要控制的是圆周积分。先除有限个零点，再让这批零点逐渐覆盖全部零点，可以避免在尚未证明可积性时直接对无穷商取边界值。

\Needspace{6\baselineskip}
\begin{lemma}[有限零点的等范数除法]\label{b3:lem:28-finite-blaschke-division}
设 \(f\in H^2(\mathbb D)\)，\(P\) 为有限 Blaschke 乘积，且 \(P\) 的每个零点在 \(f\) 中至少具有相同重数。则可去延拓后的商 \(q=f/P\) 属于 \(H^2(\mathbb D)\)，并且 \(\|q\|_{H^2}=\|f\|_{H^2}\)。空乘积或模一常数也允许。
\end{lemma}
\begin{proof}
零点的重数条件使所有商奇点可去，故 \(q\) 在整个圆盘全纯。取 \(r\) 大于 \(P\) 的全部零点的模，令
\[
 \mu_r=\min_{|z|=r}|P(z)|>0.
\]
有限乘积连续到闭圆盘且在单位圆周模为一，因此 \(\mu_r\to1\)。圆周上有
\[
 \int|q(re^{\mathrm i\theta})|^2\,\mathrm dm
 \le\mu_r^{-2}\int|f(re^{\mathrm i\theta})|^2\,\mathrm dm
 \le\mu_r^{-2}\|f\|_{H^2}^2.
\]
解析函数的平方平均由其 Taylor 系数公式可知随半径非减；令 \(r\uparrow1\)，便得到 \(q\in H^2\) 及 \(\|q\|_{H^2}\le\|f\|_{H^2}\)。反过来，由 \(f=Pq\) 和 \(|P|\le1\)，每个内部圆周上的平方平均都满足反向范数所需的不等式，故 \(\|f\|_{H^2}\le\|q\|_{H^2}\)。
\end{proof}

\Needspace{6\baselineskip}
\begin{theorem}[Hardy 函数的零点分解]\label{b3:thm:28-zero-factorization}
设 \(f\in H^2(\mathbb D)\)、\(f\not\equiv0\)。令 \(B\) 为按 \(f\) 的全部零点及重数构造的 Blaschke 乘积，原点幂与非零因子的归一化如式\eqref{b3:eq:28-blaschke-product}，整体相位任选。则
\[
 f=Bg,
 \quad g\in H^2(\mathbb D),
 \quad g(z)\ne0\quad(z\in\mathbb D),
\]
并且
\begin{equation}\label{b3:eq:28-zero-factorization-boundary}
 \begin{gathered}
 f^*=B^*g^*\quad\text{几乎处处},
 \quad |f^*|=|g^*|\quad\text{几乎处处},\\
 \|f\|_{H^2}=\|g\|_{H^2}.
 \end{gathered}
\end{equation}
\end{theorem}
\begin{proof}
\cref{b3:thm:28-hardy-zero-condition}保证相应 Blaschke 乘积存在。由于 \(B\) 与 \(f\) 具有完全相同的零点及重数，\(g=f/B\) 的商奇点全部可去，补值后在圆盘内全纯且无零点。关键尚在于证明它属于 \(H^2\)。

零点有限时，上一引理已给出结论。设非零零点列无限，并令 \(B_N\) 为包括全部原点幂的有限部分乘积。\(g_N=f/B_N\) 全纯，且上一引理给出 \(\|g_N\|_{H^2}=\|f\|_{H^2}\)。下面说明 \(g_N\to g\) 局部一致，包括共同零点处。

固定内部闭盘 \(|z|\le s\)，选 \(s<R<1\)，使圆周 \(|z|=R\) 上没有 \(B\) 的零点；零点离散，因而可避开它们的模。此圆周上的 \(|B|\) 有正下界，而 \(B_N\to B\) 一致，所以 \(g_N=f/B_N\to f/B=g\) 在该圆周上一致。最大模原理用于全纯差 \(g_N-g\)，将收敛传到 \(|z|\le R\)。于是对每个 \(0<s<1\)，
\[
 \int|g(se^{\mathrm i\theta})|^2\,\mathrm dm
 =\lim_{N\to\infty}\int|g_N(se^{\mathrm i\theta})|^2\,\mathrm dm
 \le\|f\|_{H^2}^2.
\]
这就证明 \(g\in H^2\) 及 \(\|g\|_{H^2}\le\|f\|_{H^2}\)。由 \(f=Bg\)、\(|B|\le1\) 又得反向不等式。

现在 \(B\) 有界解析且 \(g\in H^2\)，\cref{b3:prop:28-hardy-multiplier}给出 \(f^*=B^*g^*\)。\cref{b3:thm:28-blaschke-inner}给出 \(|B^*|=1\) 几乎处处，遂得到边界模相等及式\eqref{b3:eq:28-zero-factorization-boundary}。
\end{proof}

这个结论把两种作用分开：\(B\) 记录内部零点，\(g\) 保留全部边界模；但 \(B^*\) 的相位仍会改变边界函数本身。范数相等不是说函数值相同，也不是说 \(g\) 有界。整体相位的自由只会把 \(g\) 乘上一个相反的模一常数，不影响这些结论。

\subsection{零点约束与边界增长}
\label{b3:subsec:280304}

本节的必要条件也是零点列可由某个非零 \(H^2\) 函数实现的充分条件：满足 Blaschke 条件时，所构造的 \(B\) 本身就属于 \(H^2\)。因此，从“有界解析函数”放宽到“平方平均有界的解析函数”，并没有增加可实现的零点列；增加的是同一零点列可以搭配的无零点余因子。

\ProseParagraphBreak

例如取 \(\alpha>0\) 及 \(a_n=1-n^{-\alpha}\)、\(n\ge2\)。这些点能成为某个非零 \(H^2\) 函数的全部零点，当且仅当
\[
 \sum_{n=2}^{\infty}(1-a_n)=\sum_{n=2}^{\infty}n^{-\alpha}<\infty,
 \quad\text{即}\quad\alpha>1.
\]
指数越大，同样多的零点越靠近边界，每个零点在 Jensen 公式中付出的权重越小。边界聚集本身并不被禁止；受限的是带重数的总距离缺陷。

\Needspace{6\baselineskip}
\begin{corollary}[内部零点数的边界尺度]\label{b3:cor:28-zero-count-scale}
设 \(f\in H^2(\mathbb D)\)、\(f\not\equiv0\)，令 \(n_f(r)\) 为 \(|z|\le r\) 中的计重数零点个数，\(0<r<1\)。则
\[
 (1-r)\cdot n_f(r)\longrightarrow0\quad(r\uparrow1).
\]
\end{corollary}
\begin{proof}
原点的有限重数乘以 \(1-r\) 后趋于零，只需处理非零零点 \((a_j)\)。给定 \(\varepsilon>0\)，由 Blaschke 条件选 \(N\)，使 \(\sum_{j>N}(1-|a_j|)<\varepsilon\)。对 \(|a_j|\le r\) 有 \(1-r\le1-|a_j|\)，故
\[
 (1-r)\#\{\,j:|a_j|\le r\,\}
 \le N(1-r)+\sum_{j>N}(1-|a_j|)<N(1-r)+\varepsilon.
\]
先令 \(r\uparrow1\)，再令 \(\varepsilon\downarrow0\)，即得结论。
\end{proof}

与之相伴，\cref{b3:thm:28-hardy-coefficients}的 Cauchy--Schwarz 估计给出
\[
 |f(z)|\le\frac{\|f\|_{H^2}}{\sqrt{1-|z|^2}}.
\]
它是单点增长的共同上界，Blaschke 条件则是全部零点的总量约束。两者承担的任务不同。乘以任意 Blaschke 乘积可以增加许多零点，同时保持 Hardy 范数与几乎处处边界模；所以只观察边界模，不能恢复零点列。

\ProseParagraphBreak

反过来，除去全部零点后也不能认为所有边界结构都消失了。回到\secref{b3:sec:1802}中见过的无零点函数，并引入正参数 \(\tau\)：
\begin{equation}\label{b3:eq:28-singular-inner-example}
 S_\tau(z)=\exp\left(-\tau\frac{1+z}{1-z}\right),
 \quad\tau>0.
\end{equation}
由于
\[
 \operatorname{Re}\frac{1+z}{1-z}
 =\frac{1-|z|^2}{|1-z|^2}>0,
\]
它在圆盘内全纯、无零点且 \(|S_\tau|\le1\)，因此属于 \(H^2\)。对 \(e^{\mathrm i\theta}\ne1\)，公式直接连续延拓到这个边界点，并有
\[
 \frac{1+e^{\mathrm i\theta}}{1-e^{\mathrm i\theta}}
 =\mathrm i\cot\frac\theta2,
 \quad
 S_\tau(re^{\mathrm i\theta})\longrightarrow
       \exp\left(-\mathrm i\tau\cot\frac\theta2\right).
\]
极限模为一。沿径向 \(L^2\) 收敛的一条几乎处处收敛子列识别边界函数，便知 \(|S_\tau^*|=1\) 几乎处处。所以 \(S_\tau\) 是一个非恒定内函数，\(\|S_\tau\|_{H^2}=1\)；但 \(S_\tau(0)=e^{-\tau}\)，且沿实轴
\[
 S_\tau(r)=\exp\left(-\tau\frac{1+r}{1-r}\right)\longrightarrow0.
\]
一个边界点上的这种行为不影响 \(L^2\) 边界范数，却使连续到闭盘的直觉失效。

\ProseParagraphBreak

这个无零点内函数的边界奇性可以直接从对数模看见：
\[
 -\log|S_\tau(re^{\mathrm i\theta})|
 =\tau\frac{1-r^2}{|1-re^{\mathrm i\theta}|^2},
 \quad
 \int-\log|S_\tau(re^{\mathrm i\theta})|\,\mathrm dm=\tau.
\]
右端是集中于边界点 \(1\) 的质量 \(\tau\) 的 Poisson 积分。对每个避开 \(1\) 的闭圆弧，核随 \(r\uparrow1\) 一致趋于零，而它的圆周积分始终为一；负对数的质量于是集中到该点。虽然 \(-\log|S_\tau^*|=0\) 几乎处处，其内部圆周积分并不趋于零。这也解释了为什么 Blaschke 内性证明中必须先用 Jensen 公式算出积分趋零，不能仅凭几乎处处边界收敛换序。

\ProseParagraphBreak

对 \(S_\tau\) 作\cref{b3:thm:28-zero-factorization}的分解时，零点乘积可取 \(B=1\)，余因子仍是 \(S_\tau\)。更一般地，\(BS_\tau\) 与 \(B\) 有相同的全部零点和相同的边界模，却是不同的解析函数。因此本节证明的是完整的“零点分离”，并未把无零点余因子的边界结构全部分类。一般的内外因子分解还须进一步区分这种集中在零测集上的边界资料；本章后续结论不以那套分类作为未证前提。

% !TeX root = ../../Book3.tex
\section{\texorpdfstring{\(A^2\)}{A²} 的完备性与正交基}
\label{b3:sec:2804}
Hardy 范数限制每一条同心圆周上的平方平均，因而在半径趋于一时仍保留边界控制。若只要求整个圆盘上的平方积分有限，靠近边界的增长就可以更大。这个改动会把 Taylor 系数的普通平方和变成带权平方和，也会改变点值泛函的最佳常数。本节先从面积积分直接建立空间的完备性，再用圆盘的旋转对称性找出它的正交基。

\subsection{面积范数与局部控制}
\label{b3:subsec:280401}

\begin{definition}[圆盘 Bergman 空间]\label{b3:def:28-bergman-space}
设 \(\mathbb D=\{\,z\in\mathbb C\mid |z|<1\,\}\)，取归一化面积测度 \(\mathrm{d}A(z)=\mathrm{d}x\,\mathrm{d}y/\pi\)，其中 \(z=x+\mathrm{i}y\)。圆盘上的\term[bergman-kongjian]{Bergman 空间}[Bergman space]定义为
\[
 A^2(\mathbb D)\coloneqq
 \left\{\,f\in\mathcal H(\mathbb D)
 \;\middle|\;\int_{\mathbb D}|f(z)|^2\,\mathrm{d}A(z)<\infty\,\right\}.
\]
在这个空间上定义范数
\[
 \|f\|_{A^2}\coloneqq
 \left(\int_{\mathbb D}|f(z)|^2\,\mathrm{d}A(z)\right)^{1/2}.
\]
对 \(f,g\in A^2(\mathbb D)\)，规定
\[
 \langle f,g\rangle_{A^2}
 \coloneqq\int_{\mathbb D}f(z)\cdot\overline{g(z)}\,\mathrm{d}A(z).
\]
这个复内积对第一变量线性，对第二变量共轭线性。
\end{definition}
\symidx{\(A^2(\mathbb D)\)：圆盘 Bergman 空间}

归一化使 \(A(\mathbb D)=1\)，所以常数函数 \(1\) 的范数等于一。若改用未经归一化的面积 \(\mathrm{d}x\,\mathrm{d}y\)，范数将乘以 \(\sqrt\pi\)，后面的核则须除以 \(\pi\)。同一计算中的这两种约定不能混用。这里的面积范数也不同于 \(H^2\) 的边界范数，尽管它们都带上标二。

\ProseParagraphBreak

一般 \(L^2\) 函数只按几乎处处相等识别，改变单点的值不改变范数。全纯函数之间却没有这个含混：若全纯函数 \(f-g\) 在一个内点不为零，由连续性，它在某个小圆盘上的模便有正下界，因而不可能几乎处处为零。因此每个具有全纯代表的 \(L^2(\mathbb D,\mathrm{d}A)\) 等价类，至多有一个全纯代表；我们以后就使用这个代表。下面的估计进一步说明，其点值确实由面积范数连续控制。

\begin{proposition}[面积平均对点值的控制]\label{b3:prop:28-area-local-control}
设 \(f\) 在单位圆盘 \(\mathbb D\) 全纯，\(w\in\mathbb D\)，且 \(0<\rho<1-|w|\)。对于归一化面积测度 \(\mathrm{d}A=\mathrm{d}x\,\mathrm{d}y/\pi\)，有
\begin{equation}\label{b3:eq:28-area-mean-control}
 |f(w)|^2\le\frac1{\rho^2}
 \int_{D(w,\rho)}|f(z)|^2\,\mathrm{d}A(z).
\end{equation}
特别地，若 \(f\in A^2(\mathbb D)\)，则对每个 \(0\le r<1\)，
\[
 \sup_{|w|\le r}|f(w)|\le\frac{\|f\|_{A^2}}{1-r}.
\]
\end{proposition}
\begin{proof}
当 \(0<t<\rho\) 时，圆周 Cauchy 公式给出
\[
 f(w)=\frac1{2\pi}\int_0^{2\pi}
 f(w+te^{\mathrm{i}\theta})\,\mathrm{d}\theta.
\]
在总质量为一的圆周测度上使用 Cauchy--Schwarz 不等式，得到
\[
 |f(w)|^2\le\frac1{2\pi}\int_0^{2\pi}
 |f(w+te^{\mathrm{i}\theta})|^2\,\mathrm{d}\theta.
\]
两边乘以 \(2t\)，从 \(0\) 积到 \(\rho\)。左边为 \(\rho^2|f(w)|^2\)，右边由极坐标换元恰好等于归一化面积积分。这证明了\eqref{b3:eq:28-area-mean-control}。

若 \(|w|\le r\)，可以对全部这些 \(w\) 使用同一个 \(0<\rho<1-r\)，并把小圆盘积分放大为整个 \(\mathbb D\) 上的积分。于是 \(\sup_{|w|\le r}|f(w)|\le\|f\|_{A^2}/\rho\)。令 \(\rho\uparrow1-r\) 即得结论；无需假设 \(f\) 连续延拓到单位圆周。
\end{proof}

这个论证是\subsecref{b3:subsec:040302}所用局部平方平均估计的圆盘版本，这里已经按 \(\mathrm{d}A=\mathrm{d}x\,\mathrm{d}y/\pi\) 重写常数。它说明面积积分虽不单独限制一般可测函数的点值，却能限制全纯函数在每个固定内部紧集上的大小。全纯性通过圆周平均，把一个内点与周围一片正面积区域联系起来。

\ProseParagraphBreak

靠近边界时，上界允许增长。这个允许量有实际内容：\(\sqrt{n+1}\,z^n\) 的面积范数稍后将算出恒为一，但它在接近单位圆周处可以很大。在固定的 \(|z|\le r<1\) 上，\(\sqrt{n+1}\,r^n\to0\)，又与局部控制完全相容。下一节将把当前的粗常数改进到最佳值。

\subsection{Bergman 空间的完备性}
\label{b3:subsec:280402}

\(L^2\) 的完备性首先给出可测函数类的极限；要建立 \(A^2\) 的完备性，还须把这一极限识别为全纯函数。局部估计正好连接这两个空间。它将面积范数中的 Cauchy 条件变为每个闭子圆盘上的一致 Cauchy 条件；得到全纯极限以后，再回到面积积分证明范数收敛。

\begin{theorem}[圆盘 Bergman 空间的完备性]\label{b3:thm:28-bergman-complete}
在单位圆盘 \(\mathbb D\) 上取 \(\mathrm{d}A=\mathrm{d}x\,\mathrm{d}y/\pi\)。空间 \(A^2(\mathbb D)\) 关于内积 \(\langle f,g\rangle_{A^2}=\int_{\mathbb D}f\overline g\,\mathrm{d}A\) 是复 Hilbert 空间。以唯一的全纯代表将它视为 \(L^2(\mathbb D,\mathrm{d}A)\) 的子空间时，这个子空间是闭的。
\end{theorem}
\begin{proof}
线性组合仍全纯且平方可积，内积的各项性质来自积分；正定性由全纯代表的唯一性保证。只需证明完备性与闭性。

设 \((f_j)\) 是 \(A^2\) 中的 Cauchy 列。对固定 \(0<r<1\)，由\cref{b3:prop:28-area-local-control}作用于 \(f_j-f_k\)，有
\[
 \sup_{|z|\le r}|f_j(z)-f_k(z)|
 \le\frac1{1-r}\|f_j-f_k\|_{A^2}\longrightarrow0
 \quad(j,k\to\infty).
\]
所以每个闭子圆盘上都有一致极限。在两个闭子圆盘的重叠处，这些极限等于同一数列的逐点极限，因而彼此相容，给出 \(\mathbb D\) 上的函数 \(f\)。任何内部紧集都包含于某个闭子圆盘，故 \(f_j\to f\) 局部一致。各 \(f_j\) 全纯，\cref{b3:prop:04-holomorphic-limit}于是保证 \(f\) 全纯。

还须确认 \(f\) 具有有限的全盘面积范数，而不仅在每个紧集上良好。Cauchy 列有界，设 \(\|f_j\|_{A^2}\le M\)。对非负函数 \(|f_j|^2\) 使用 Fatou 引理，得到
\[
 \int_{\mathbb D}|f|^2\,\mathrm{d}A
 \le\liminf_{j\to\infty}\int_{\mathbb D}|f_j|^2\,\mathrm{d}A
 \le M^2,
\]
故 \(f\in A^2\)。给定 \(\varepsilon>0\)，选择 \(N\) 使 \(j,k\ge N\) 时 \(\|f_j-f_k\|_{A^2}<\varepsilon\)。固定 \(j\ge N\)，再让 \(k\to\infty\)，对 \(|f_j-f_k|^2\) 使用 Fatou 引理，得
\[
 \|f_j-f\|_{A^2}^2
 \le\liminf_{k\to\infty}\|f_j-f_k\|_{A^2}^2
 \le\varepsilon^2.
\]
因此整列在面积范数中趋于 \(f\)，完备性得证。

最后，若 \(f_j\in A^2\) 在 \(L^2(\mathbb D,\mathrm{d}A)\) 中趋于某个等价类 \(u\)，则它也是上述 Cauchy 列，故存在全纯 \(f\in A^2\) 使 \(\|f_j-f\|_{L^2}\to0\)。三角不等式给出 \(\|u-f\|_{L^2}=0\)，所以 \(u\) 正是这个全纯代表的等价类。这证明了子空间的闭性。
\end{proof}

这里两次使用 Fatou 引理的目的不同：第一次证明候选极限仍平方可积，第二次把 Cauchy 误差传给最终极限。我们没有从局部一致收敛直接断言全盘积分收敛，因为边界附近仍可能有质量集中。例如 \(\sqrt{n+1}\,z^n\) 局部一致趋于零，面积范数却不趋于零；它不是面积范数中的 Cauchy 列。

\ProseParagraphBreak

范数收敛还控制固定阶导数。若 \(f_j\to f\) 于 \(A^2\)，取 \(0<r<s<1\)，先在 \(|z|\le s\) 上使用局部估计，再在每个 \(|w|\le r\) 周围使用半径 \(s-r\) 的 Cauchy 导数估计，便有
\[
 \sup_{|w|\le r}|f_j^{(m)}(w)-f^{(m)}(w)|
 \le\frac{m!}{(s-r)^m(1-s)}\|f_j-f\|_{A^2}
 \quad(m=0,1,2,\ldots).
\]
这使面积范数中的极限与前文熟悉的全纯函数局部极限相容，但没有把任何结论延伸到移动的边界点上。

\subsection{单项式的正交性}
\label{b3:subsec:280403}

圆盘与面积测度都在旋转下不变。因此不同角频率应当互相正交；单项式 \(z^n=r^ne^{\mathrm{i}n\theta}\) 正好把角频率与径向幂次分开。直接计算还能看出，为什么面积空间的归一化基与 Hardy 空间不同。

\begin{proposition}[单项式的面积内积]\label{b3:prop:28-monomial-inner-product}
在 \(A^2(\mathbb D)\) 中取 \(\mathrm{d}A=\mathrm{d}x\,\mathrm{d}y/\pi\)。对任意非负整数 \(n,m\)，有
\[
 \langle z^n,z^m\rangle_{A^2}
 =\begin{cases}
 0,&n\ne m,\\
 \dfrac1{n+1},&n=m.
 \end{cases}
\]
因此函数 \(e_n(z)=\sqrt{n+1}\,z^n\)，\(n=0,1,2,\ldots\)，构成 \(A^2(\mathbb D)\) 的标准正交系。
\end{proposition}
\begin{proof}
记 \(\mathrm{d}m(\theta)=\mathrm{d}\theta/(2\pi)\)。极坐标面积元为 \(\mathrm{d}A=2r\,\mathrm{d}r\,\mathrm{d}m\)，所以
\[
 \langle z^n,z^m\rangle_{A^2}
 =2\int_0^1r^{n+m+1}\,\mathrm{d}r
 \int_0^{2\pi}e^{\mathrm{i}(n-m)\theta}\,\mathrm{d}m(\theta).
\]
当 \(n\ne m\) 时，角积分为零；当 \(n=m\) 时，角积分为一，而径向积分为 \(2/(2n+2)=1/(n+1)\)。乘上 \(\sqrt{n+1}\) 正好把每个单项式的长度变成一。
\end{proof}

对于多项式 \(p(z)=\sum_{n=0}^Na_nz^n\)，有限和的正交性立即给出
\[
 \|p\|_{A^2}^2=\sum_{n=0}^N\frac{|a_n|^2}{n+1}.
\]
例如 \(\|1+2z\|_{A^2}^2=1+4/2=3\)，而它的 Hardy 范数平方为 \(1+4=5\)。面积积分对高次项给出的权重较小，是因为 \(r^n\) 只有靠近边界时才不小；把全部半径再积一次，会减轻这一圈的贡献。

\ProseParagraphBreak

这些单项式相互正交，还需要知道它们能否逼近任意 \(A^2\) 函数。Taylor 展开与非负项积分将说明，多项式逼近可以从内部圆盘上的一致收敛提升为整个圆盘上的面积范数收敛。

\subsection{系数表示与正交基的完备性}
\label{b3:subsec:280404}

\begin{theorem}[Bergman 范数与 Taylor 系数]\label{b3:thm:28-bergman-coefficients}
设 \(f(z)=\sum_{n=0}^{\infty}a_nz^n\) 在单位圆盘 \(\mathbb D\) 全纯，并取 \(\mathrm{d}A=\mathrm{d}x\,\mathrm{d}y/\pi\)。则
\begin{equation}\label{b3:eq:28-bergman-coefficient-norm}
 \int_{\mathbb D}|f(z)|^2\,\mathrm{d}A(z)
 =\sum_{n=0}^{\infty}\frac{|a_n|^2}{n+1},
\end{equation}
允许两边同为 \(+\infty\)。特别地，\(f\in A^2(\mathbb D)\) 当且仅当右端有限；在此情形，Taylor 部分和 \(p_N(z)=\sum_{n=0}^Na_nz^n\) 在 \(A^2\) 范数中趋于 \(f\)，且
\[
 \|f-p_N\|_{A^2}^2=\sum_{n>N}\frac{|a_n|^2}{n+1}.
\]
\end{theorem}
\begin{proof}
固定 \(0<r<1\)。Taylor 级数在圆周 \(|z|=r\) 上一致收敛，因此有限部分和的平方模积分趋于 \(f\) 的平方模积分。先对有限和使用角频率正交性，再取极限，得到
\[
 \int_0^{2\pi}|f(re^{\mathrm{i}\theta})|^2\,\mathrm{d}m(\theta)
 =\sum_{n=0}^{\infty}|a_n|^2r^{2n}.
\]


对于 \(0<R<1\)，两边乘以 \(2r\)，从 \(0\) 积到 \(R\)。右端各项非负，Tonelli 定理允许交换求和与径向积分，给出
\[
 \int_{|z|<R}|f(z)|^2\,\mathrm{d}A(z)
 =\sum_{n=0}^{\infty}\frac{|a_n|^2}{n+1}R^{2n+2}.
\]
令 \(R\uparrow1\)，左边由区域递增的单调收敛趋于全盘积分；右边也由非负项的单调收敛趋于所述级数。这证明了\eqref{b3:eq:28-bergman-coefficient-norm}，包括无穷值情形。将同一公式用于全纯函数 \(f-p_N\)，便得到精确的尾范数；在级数有限时，这个尾和趋于零。
\end{proof}

反方向也需要说明：任意满足 \(\sum_{n\ge0}|a_n|^2/(n+1)<\infty\) 的复数列，都确实给出圆盘上的全纯函数。对 \(0\le r<1\)，Cauchy--Schwarz 不等式给出
\begin{equation}\label{b3:eq:28-weighted-series-control}
 \begin{aligned}
 \sum_{n=0}^{\infty}|a_n|r^n
 &\le\left(\sum_{n=0}^{\infty}\frac{|a_n|^2}{n+1}\right)^{1/2}
 \left(\sum_{n=0}^{\infty}(n+1)r^{2n}\right)^{1/2}
 \\
 &=\frac1{1-r^2}
 \left(\sum_{n=0}^{\infty}\frac{|a_n|^2}{n+1}\right)^{1/2}.
 \end{aligned}
\end{equation}
最后的求和由几何级数逐项求导得到。对尾级数使用同一估计，可知幂级数在每个闭子圆盘上一致绝对收敛，因而定义全纯函数；再用\cref{b3:thm:28-bergman-coefficients}即可确认它属于 \(A^2\)。

\begin{corollary}[Bergman 标准正交基]\label{b3:cor:28-bergman-basis}
在单位圆盘的归一化面积 Bergman 空间 \(A^2(\mathbb D)\) 中，函数
\[
 e_n(z)=\sqrt{n+1}\,z^n\quad(n=0,1,2,\ldots)
\]
构成标准正交基。若 \(f(z)=\sum_{n\ge0}a_nz^n\in A^2(\mathbb D)\)，则
\[
 \langle f,e_n\rangle_{A^2}=\frac{a_n}{\sqrt{n+1}},
 \quad
 f=\sum_{n=0}^{\infty}\langle f,e_n\rangle_{A^2}e_n,
\]
后一个级数同时在面积范数及局部一致意义下收敛。映射
\[
 f\longmapsto\left(\frac{a_n}{\sqrt{n+1}}\right)_{n\ge0}
\]
是从 \(A^2(\mathbb D)\) 到复序列空间 \(\ell^2(\mathbb N_0)\) 的满等距线性映射。
\end{corollary}
\begin{proof}
Taylor 部分和的面积范数收敛表明，多项式的线性空间在 \(A^2\) 中稠密，故先前的标准正交系是基。对部分和与 \(e_n\) 作内积，有限正交计算给出 \(a_n/\sqrt{n+1}\)；内积连续性允许取部分和的范数极限，得到同一公式。范数恒等式给出等距性；对于任意 \((c_n)\in\ell^2\)，取 \(a_n=\sqrt{n+1}\,c_n\)，前述加权级数构造给出相应的 \(A^2\) 函数，因此映射是满射。局部一致收敛也可由面积范数收敛和\cref{b3:prop:28-area-local-control}直接推出。
\end{proof}

现在每个全纯平方可积函数都由一组平方可和的正交坐标表示。与 Hardy 空间的差别集中在权重 \(1/(n+1)\)：它既说明面积可积容许较大的 Taylor 系数，也决定下一节的核中为何出现二次分母。

% !TeX root = ../../Book3.tex
\section{圆盘上的 Bergman 核}
\label{b3:sec:2805}
在面积平方可积的全纯函数中，知道范数能把某个内点的函数值限制到什么程度？上一节给出了一个局部平均上界，正交坐标则能进一步找到最佳上界及取得等号的函数。表示点值的核还会给出一个积分算子：它从任意面积平方可积函数中提取全纯部分。这里依次构造核、计算核、建立投影，最后与 Hardy 空间比较。

\subsection{点值泛函与核的存在性}
\label{b3:subsec:280501}

固定 \(w\in\mathbb D\)。映射 \(f\mapsto f(w)\) 对 \(f\) 线性，且由\cref{b3:prop:28-area-local-control}连续。我们希望找到某个 \(K_w\in A^2(\mathbb D)\)，使
\[
 f(w)=\langle f,K_w\rangle_{A^2}
 \quad(f\in A^2(\mathbb D)).
\]
这就是前面 Hardy 核所具有的再生性质。构造时不必先引入一般 Hilbert 空间上的泛函表示定理：上一节的标准正交基已经给出所需向量的全部坐标。

\ProseParagraphBreak

具体地，写 \(e_n(z)=\sqrt{n+1}\,z^n\)，并设 \(f=\sum_{n\ge0}c_ne_n\)。我们需要内积右端等于 \(\sum_{n\ge0}c_n\cdot e_n(w)\)。因为内积对第二变量共轭线性，核的第 \(n\) 个坐标应取 \(\overline{e_n(w)}\)，而不是 \(e_n(w)\)。这一共轭正是二元核对第一个变量全纯、对第二个变量反全纯的来源。

\begin{proposition}[Bergman 核的正交构造]\label{b3:prop:28-bergman-kernel-construction}
在归一化面积测度 \(\mathrm{d}A=\mathrm{d}x\,\mathrm{d}y/\pi\) 下，设 \(A^2(\mathbb D)\) 的内积对第一变量线性，且 \(e_n(z)=\sqrt{n+1}\,z^n\)。对每个 \(w\in\mathbb D\)，级数
\begin{equation}\label{b3:eq:28-bergman-kernel-series}
 K_w\coloneqq\sum_{n=0}^{\infty}\overline{e_n(w)}e_n
\end{equation}
在 \(A^2(\mathbb D)\) 中收敛，并且
\begin{equation}\label{b3:eq:28-bergman-reproduction}
 f(w)=\langle f,K_w\rangle_{A^2}
 \quad(f\in A^2(\mathbb D)).
\end{equation}
具有这一性质的 \(K_w\) 唯一。二元函数 \(K(z,w)\coloneqq K_w(z)\) 称为这个空间的\term[bergman-he]{Bergman 核}[Bergman kernel]。
\end{proposition}
\symidx{\(K(z,w)=K_w(z)\)：归一化面积圆盘 Bergman 核}
\begin{proof}
对固定 \(|w|<1\)，有
\[
 \sum_{n=0}^{\infty}|e_n(w)|^2
 =\sum_{n=0}^{\infty}(n+1)|w|^{2n}
 =\frac1{(1-|w|^2)^2}<\infty.
\]
因此\eqref{b3:eq:28-bergman-kernel-series}的部分和在面积范数中构成 Cauchy 列；\cref{b3:thm:28-bergman-complete}给出 \(K_w\in A^2\)。对于 \(f=\sum c_ne_n\)，内积的连续性和正交性给出
\[
 \langle f,K_w\rangle_{A^2}
 =\sum_{n=0}^{\infty}c_n\cdot e_n(w)=f(w).
\]
标量级数绝对收敛，因为 \((c_n)\) 与 \((e_n(w))\) 都平方可和；最后一个等号也与\cref{b3:cor:28-bergman-basis}的局部一致展开相符。

如果另有 \(L_w\in A^2\) 对全部 \(f\) 满足相同再生公式，则 \(\langle f,K_w-L_w\rangle_{A^2}=0\)。取 \(f=K_w-L_w\)，得到 \(\|K_w-L_w\|_{A^2}^2=0\)，所以两个全纯函数处处相等。
\end{proof}

核不仅是求值公式的记号。把 \(f=K_w\) 代入再生公式，可得 \(K(w,w)=\|K_w\|_{A^2}^2>0\)；把 \(f=K_v\) 代入，其中 \(v\in\mathbb D\)，又得
\[
 \langle K_v,K_w\rangle_{A^2}=K_v(w)=K(w,v),
 \quad K(z,w)=\overline{K(w,z)}.
\]
因此核在两个点处的值，还记录了两个点值表示向量之间的内积。第二个等式也可以直接从正交级数逐项取共轭得到。

\subsection{圆盘 Bergman 核的计算}
\label{b3:subsec:280502}

几何级数在 \(|t|<1\) 内满足 \(\sum_{n\ge0}t^n=1/(1-t)\)。对 \(t/(1-t)=\sum_{n\ge0}t^{n+1}\) 求导，便得
\[
 \sum_{n=0}^{\infty}(n+1)t^n=\frac1{(1-t)^2}.
\]
在核的正交展开中代入 \(t=z\overline w\)，由于 \(|z|<1\)、\(|w|<1\)，收敛条件确实满足。

\begin{theorem}[圆盘 Bergman 核与最佳点值界]\label{b3:thm:28-bergman-kernel}
在 \(\mathbb D\) 上取归一化面积 \(\mathrm{d}A=\mathrm{d}x\,\mathrm{d}y/\pi\)。则 \(A^2(\mathbb D)\) 的再生核为
\begin{equation}\label{b3:eq:28-bergman-kernel-formula}
 K_w(z)=K(z,w)=\frac1{(1-z\overline w)^2}
 \quad(z,w\in\mathbb D).
\end{equation}
对每个 \(f\in A^2(\mathbb D)\) 和每个 \(w\in\mathbb D\)，有
\begin{equation}\label{b3:eq:28-bergman-evaluation}
 \begin{aligned}
 f(w)=\int_{\mathbb D}\frac{f(z)}{(1-w\overline z)^2}\,\mathrm{d}A(z),
 \\
 |f(w)|\le\frac{\|f\|_{A^2}}{1-|w|^2}.
 \end{aligned}
\end{equation}
点值泛函的范数恰为 \(1/(1-|w|^2)\)。不等式取等号，当且仅当 \(f\) 是 \(K_w\) 的复数倍，零函数也包括在内。
\end{theorem}
\begin{proof}
正交级数与上述几何求和给出核的显式公式。将它代入\eqref{b3:eq:28-bergman-reproduction}，注意内积第二变量要取共轭，便得到积分表达式。

由正交展开或再生公式，
\[
 \|K_w\|_{A^2}^2=K(w,w)=\frac1{(1-|w|^2)^2}.
\]
Cauchy--Schwarz 不等式于是给出
\[
 |f(w)|=|\langle f,K_w\rangle_{A^2}|
 \le\|f\|_{A^2}\cdot\|K_w\|_{A^2}
 =\frac{\|f\|_{A^2}}{1-|w|^2}.
\]
积分中的乘积也因这同一个不等式绝对可积。取单位向量 \(K_w/\|K_w\|_{A^2}\) 时取到上界，故常数不能再减小。由于 \(K_w\ne0\)，Cauchy--Schwarz 的等号条件恰为 \(f\) 与 \(K_w\) 线性相关。
\end{proof}

在 \(w=0\) 时，核为常数一，再生公式成为 \(f(0)=\int_{\mathbb D}f\,\mathrm{d}A\)，是面积平均公式。一般内点的核会向相应的边界方向集中，最佳点值上界则由原点处的一增大到 \(1/(1-|w|^2)\)。除原点处两个上界相同外，它严格改善了上一节的粗上界 \(1/(1-|w|)\)，但仍允许靠近单位圆周时增长。

\begin{corollary}[指定一点值的最小面积范数]\label{b3:cor:28-bergman-extremal}
设 \(w\in\mathbb D\)、\(a\in\mathbb C\)，面积测度仍为 \(\mathrm{d}A=\mathrm{d}x\,\mathrm{d}y/\pi\)。在全部满足 \(f\in A^2(\mathbb D)\) 且 \(f(w)=a\) 的函数中，范数的最小值为 \(|a|(1-|w|^2)\)，唯一的极小元为
\[
 f_{w,a}(z)=a\,\frac{(1-|w|^2)^2}{(1-z\overline w)^2}.
\]
对任意这样的 \(f\)，还成立
\[
 \|f\|_{A^2}^2
 =|a|^2(1-|w|^2)^2+\|f-f_{w,a}\|_{A^2}^2.
\]
\end{corollary}
\begin{proof}
令 \(f_{w,a}=aK_w/K(w,w)\)，则 \(f_{w,a}(w)=a\)。函数 \(h=f-f_{w,a}\) 满足 \(h(w)=0\)，再生公式给出 \(\langle h,K_w\rangle_{A^2}=0\)。因此 \(h\) 与 \(f_{w,a}\) 正交，展开平方范数即得所述恒等式；其中
\[
 \|f_{w,a}\|_{A^2}^2
 =\frac{|a|^2}{K(w,w)}=|a|^2(1-|w|^2)^2.
\]
右边剩余平方范数非负，且仅在 \(h=0\) 时为零，证明了极小值及唯一性。这一论证也涵盖 \(a=0\)。
\end{proof}

极值函数的形状还可从圆盘自同构理解。对 \(a\in\mathbb D\)，令
\[
 \varphi_a(z)=\frac{a-z}{1-\overline a z},
 \quad
 \varphi_a'(z)=-\frac{1-|a|^2}{(1-\overline a z)^2}.
\]
该映射交换 \(a\) 与原点，且 \(\varphi_a\circ\varphi_a\) 为恒等映射。直接通分给出
\begin{equation}\label{b3:eq:28-disc-automorphism-identity}
 1-\varphi_a(z)\cdot\overline{\varphi_a(w)}
 =\frac{(1-|a|^2)(1-z\overline w)}
 {(1-\overline a z)(1-a\overline w)}.
\end{equation}
取 \(w=z\) 时，右端为正，故 \(\varphi_a\) 确把圆盘映入自身；配合其自逆性，便知它是圆盘的双全纯自同构。

\begin{proposition}[自同构下的面积范数与核]\label{b3:prop:28-bergman-automorphism}
设 \(a\in\mathbb D\)、\(\theta\in\mathbb R\)，并令 \(\varphi(z)=e^{\mathrm{i}\theta}(a-z)/(1-\overline a z)\)。在归一化面积 Bergman 空间 \(A^2(\mathbb D)\) 上，算子
\[
 (U_\varphi f)(z)\coloneqq\varphi'(z)\cdot f(\varphi(z))
\]
是满等距线性映射，其逆为 \(U_{\varphi^{-1}}\)。Bergman 核满足
\begin{equation}\label{b3:eq:28-bergman-transformation}
 K(z,w)=\varphi'(z)\cdot\overline{\varphi'(w)}
 \cdot K(\varphi(z),\varphi(w))
 \quad(z,w\in\mathbb D).
\end{equation}
\end{proposition}
\begin{proof}
复导数非零的全纯映射作为实二维映射的 Jacobi 行列式等于 \(|\varphi'|^2\)：若 \(\varphi'=u+\mathrm{i}v\)，其微分矩阵为 \(\left(\begin{smallmatrix}u&-v\\v&u\end{smallmatrix}\right)\)，行列式为 \(u^2+v^2\)。因此换元 \(\zeta=\varphi(z)\) 给出
\[
 \int_{\mathbb D}|\varphi'(z)|^2\cdot|f(\varphi(z))|^2\,\mathrm{d}A(z)
 =\int_{\mathbb D}|f(\zeta)|^2\,\mathrm{d}A(\zeta).
\]
这证明算子保范数，也证明它把 \(A^2\) 映入自身。链式法则给出 \(U_\varphi U_{\varphi^{-1}}=U_{\varphi^{-1}}U_\varphi=I\)，从而得到满射性及逆算子。

当 \(\varphi=\varphi_a\) 时，将\eqref{b3:eq:28-disc-automorphism-identity}和导数公式代入核的显式表达，分子分母相消即得\eqref{b3:eq:28-bergman-transformation}。再乘一个旋转 \(e^{\mathrm{i}\theta}\)，既不改变 \(\varphi(z)\cdot\overline{\varphi(w)}\)，也不改变两个导数的乘积 \(\varphi'(z)\cdot\overline{\varphi'(w)}\)，所以等式对所述全部参数成立。
\end{proof}

面积范数在变量替换时需要导数因子，因为二维面积按 \(|\varphi'|^2\) 伸缩。这个因子也解释了核的二次分母：将原点处的常数极值函数通过 \(U_{\varphi_a}\) 搬到另一点，会出现 \(\varphi_a'\) 所含的 \((1-\overline a z)^{-2}\)。单纯把函数与 \(\varphi_a\) 复合，并不保持当前的面积范数。

\subsection{正交展开与 Bergman 投影}
\label{b3:subsec:280503}

再生公式的左边只对全纯函数定义，但右边的积分对任何 \(u\in L^2(\mathbb D,\mathrm{d}A)\) 都有意义。这提示我们用同一个核将一般函数变成全纯函数。要证明它确实是正交投影，先从正交系建立算子，再辨认它的积分表达，能把收敛问题处理得更清楚。

\begin{theorem}[Bergman 投影]\label{b3:thm:28-bergman-projection}
设 \(\mathbb D\) 为单位圆盘，\(\mathrm{d}A=\mathrm{d}x\,\mathrm{d}y/\pi\)，并在复空间 \(L^2(\mathbb D,\mathrm{d}A)\) 上取第一变量线性的内积。对 \(u\in L^2(\mathbb D,\mathrm{d}A)\)，定义
\begin{equation}\label{b3:eq:28-bergman-projection-series}
 P_Au\coloneqq\sum_{n=0}^{\infty}\langle u,e_n\rangle_{L^2}e_n,
 \quad e_n(z)=\sqrt{n+1}\,z^n.
\end{equation}
这个级数在 \(A^2(\mathbb D)\) 范数中收敛，并确定从 \(L^2(\mathbb D,\mathrm{d}A)\) 到其闭子空间 \(A^2(\mathbb D)\) 的正交投影，称为\term[bergman-touying]{Bergman 投影}[Bergman projection]。它满足
\[
 P_A^2=P_A,
 \quad \|P_Au\|_{A^2}\le\|u\|_{L^2},
 \quad \|P_A\|=1,
\]
并有对每个 \(w\in\mathbb D\) 成立的点态积分表达式
\begin{equation}\label{b3:eq:28-bergman-projection-integral}
 (P_Au)(w)=\int_{\mathbb D}\frac{u(z)}{(1-w\overline z)^2}\,\mathrm{d}A(z).
\end{equation}
\end{theorem}
\symidx{\(P_A\)：从面积 \(L^2\) 到 \(A^2(\mathbb D)\) 的 Bergman 投影}
\begin{proof}
记 \(c_n=\langle u,e_n\rangle_{L^2}\)，并令 \(q_N=\sum_{n=0}^Nc_ne_n\)。有限正交展开给出
\[
 \|u-q_N\|_{L^2}^2
 =\|u\|_{L^2}^2-\sum_{n=0}^N|c_n|^2\ge0.
\]
所以 \(\sum_{n\ge0}|c_n|^2\le\|u\|_{L^2}^2\)。这就是当前系统所需的 Bessel 不等式，也说明 \((q_N)\) 在 \(A^2\) 中为 Cauchy 列。由其完备性，存在 \(P_Au\in A^2\) 使 \(q_N\to P_Au\) 于面积范数，且
\[
 \|P_Au\|_{A^2}^2=\sum_{n=0}^{\infty}|c_n|^2\le\|u\|_{L^2}^2.
\]
系数逐项线性以及极限唯一性还证明 \(P_A\) 线性。

对每个 \(n\)，有 \(\langle u-P_Au,e_n\rangle_{L^2}=0\)。任取 \(h\in A^2\)，以它的 Taylor 多项式在面积范数中逼近，并用内积连续性取极限，便得
\[
 \langle u-P_Au,h\rangle_{L^2}=0.
\]
于是 \(u-P_Au\) 与整个 \(A^2\) 正交。当 \(u\in A^2\) 时，\cref{b3:cor:28-bergman-basis}给出 \(P_Au=u\)，故 \(P_A^2=P_A\)。上面的范数界说明 \(\|P_A\|\le1\)，而 \(P_A1=1\)、\(\|1\|_{L^2}=1\)，给出 \(\|P_A\|\ge1\)。

最后固定 \(w\in\mathbb D\)。核的部分和
\[
 K_{w,N}(z)=\sum_{n=0}^N\overline{e_n(w)}\cdot e_n(z)
\]
在 \(L^2\) 中趋于 \(K_w\)。另一方面，\(q_N\to P_Au\) 于 \(A^2\)，所以局部控制保证 \(q_N(w)\to(P_Au)(w)\)。由有限和恒等式
\[
 q_N(w)=\langle u,K_{w,N}\rangle_{L^2}
\]
及 Cauchy--Schwarz 估计
\[
 |\langle u,K_{w,N}-K_w\rangle_{L^2}|
 \le\|u\|_{L^2}\cdot\|K_{w,N}-K_w\|_{L^2}\longrightarrow0,
\]
得到 \((P_Au)(w)=\langle u,K_w\rangle_{L^2}\)，即\eqref{b3:eq:28-bergman-projection-integral}。同一不等式也保证该积分绝对收敛。
\end{proof}

这里通过有限和与 \(L^2\) 极限得到积分公式，没有把任意形式的无穷级数直接搬进积分号。若希望直接检查换序，对固定 \(|w|\le r<1\)，也有
\[
 \sum_{n=0}^{\infty}(n+1)|w|^n\cdot|z|^n
 \le\sum_{n=0}^{\infty}(n+1)r^n=\frac1{(1-r)^2}.
\]
由于 \(A(\mathbb D)=1\)，Cauchy--Schwarz 给出 \(\|u\|_{L^1}\le\|u\|_{L^2}\)，故 \(|u(z)|/(1-r)^2\) 是共同可积控制。支配收敛或绝对可积的 Fubini 换序因此可用，而且只保证在固定内部闭圆盘上的一致控制。

\ProseParagraphBreak

正交性使 \(P_Au\) 具有最佳逼近意义。对于任意 \(h\in A^2\)，展开
\[
 u-h=(u-P_Au)+(P_Au-h)
\]
并利用两项正交，得到
\[
 \|u-h\|_{L^2}^2
 =\|u-P_Au\|_{L^2}^2+\|P_Au-h\|_{A^2}^2.
\]
因此 \(P_Au\) 是距离 \(u\) 最近的唯一全纯平方可积函数。它不是对原函数作逐点修补，而是同时利用全盘积分，在面积平方误差中找出最优的全纯近似。

\ProseParagraphBreak

混合单项式可以直接看清这个过程。设 \(p,q\) 为非负整数，取 \(u(z)=z^p\overline z^{\,q}\)。当 \(n\ge0\) 时，
\[
 \langle u,e_n\rangle_{L^2}
 =\sqrt{n+1}\int_{\mathbb D}z^p\overline z^{\,q+n}\,\mathrm{d}A(z).
\]
角积分只在 \(p=q+n\) 时非零。若 \(p\ge q\)，唯一留下的指标为 \(n=p-q\)，径向积分为 \(1/(p+1)\)；若 \(p<q\)，则没有非负指标可留下。因此
\begin{equation}\label{b3:eq:28-projection-mixed-monomial}
 P_A(z^p\overline z^{\,q})
 =\begin{cases}
 \dfrac{p-q+1}{p+1}\,z^{p-q},&p\ge q,\\
 0,&p<q.
 \end{cases}
\end{equation}
例如 \(P_A(\overline z)=0\)，而 \(P_A(|z|^2)=1/2\)、\(P_A(z^2\overline z)=2z/3\)。尤其不能把“投影到全纯函数”理解为在表达式中简单删掉每个带 \(\overline z\) 的项；\(|z|^2\) 虽非全纯，却有非零的常数投影。

\subsection{Hardy 核与 Bergman 核的比较}
\label{b3:subsec:280504}

两个空间都由全纯函数组成，都有单项式基，也都能通过内积再生点值。真正不同的是积分所在的位置。Hardy 空间控制所有径向圆周平均，得到单位圆周上的 \(L^2\) 边界函数；Bergman 空间把这些平均再按 \(2r\,\mathrm{d}r\) 积起来，因而允许平均值在 \(r\uparrow1\) 时无界。\cref{b3:tab:28-hardy-bergman}把已证明的结论放在同一测度约定下比较。

\begin{table}[htbp]
\centering
\caption[圆盘 Hardy 与 Bergman 空间]{圆盘 Hardy 与 Bergman 空间。这里 \(f(z)=\sum_{n\ge0}a_nz^n\)，\(\mathrm{d}m=\mathrm{d}\theta/(2\pi)\)，\(\mathrm{d}A=\mathrm{d}x\,\mathrm{d}y/\pi\)；边界收敛一行专指本章已证明的 \(L^2\) 收敛。}
\label{b3:tab:28-hardy-bergman}
\small
\begin{tabularx}{\textwidth}{@{}p{5.5em}XX@{}}
\toprule
比较项目 & \(H^2(\mathbb D)\) & \(A^2(\mathbb D)\)\\
\midrule
范数的测度 & 圆周 \(\mathrm{d}m\)，对半径取上确界 & 圆盘面积 \(\mathrm{d}A\)\\
系数范数平方 & \(\sum_{n\ge0}|a_n|^2\) & \(\sum_{n\ge0}|a_n|^2/(n+1)\)\\
标准正交基 & \(z^n\)，\(n\ge0\) & \(\sqrt{n+1}\,z^n\)，\(n\ge0\)\\
径向边界 & 有唯一 \(L^2(\mathrm{d}m)\) 边界极限，且保范数 & 面积范数有限不保证有 \(L^2(\mathrm{d}m)\) 边界极限\\
再生核 & \(k_w(z)=(1-z\overline w)^{-1}\) & \(K_w(z)=(1-z\overline w)^{-2}\)\\
点值泛函范数 & \((1-|w|^2)^{-1/2}\) & \((1-|w|^2)^{-1}\)\\
\bottomrule
\end{tabularx}
\end{table}

表中 Hardy 部分分别来自\cref{b3:thm:28-hardy-coefficients}、\cref{b3:thm:28-hardy-boundary}与\cref{b3:thm:28-hardy-kernel}。边界结论说的是 \(f(re^{\mathrm{i}\theta})\) 在 \(L^2(\mathrm{d}m)\) 中的径向极限；不能仅凭这个范数收敛，把对全部半径的几乎处处逐点极限另当成已经证明的结论。Bergman 核的积分发生在圆盘内部，也不需要先赋予一般 \(A^2\) 函数一个平方可积的边界函数。

\begin{proposition}[两个空间的包含关系]\label{b3:prop:28-hardy-bergman-inclusion}
在单位圆盘 \(\mathbb D\) 上，Hardy 空间使用归一化圆周测度 \(\mathrm{d}m=\mathrm{d}\theta/(2\pi)\)，Bergman 空间使用归一化面积 \(\mathrm{d}A=\mathrm{d}x\,\mathrm{d}y/\pi\)。恒等映射给出连续、稠密且严格的包含
\[
 H^2(\mathbb D)\subsetneq A^2(\mathbb D),
 \quad \|f\|_{A^2}\le\|f\|_{H^2}\quad(f\in H^2(\mathbb D)).
\]
该包含映射的范数为一，而 \(H^2(\mathbb D)\) 在 \(A^2(\mathbb D)\) 的面积范数中不是闭子空间。
\end{proposition}
\begin{proof}
对于 \(f(z)=\sum a_nz^n\in H^2\)，系数范数公式给出
\[
 \|f\|_{A^2}^2=\sum_{n=0}^{\infty}\frac{|a_n|^2}{n+1}
 \le\sum_{n=0}^{\infty}|a_n|^2=\|f\|_{H^2}^2.
\]
常数函数一使不等式取等号，故映射范数为一。所有多项式都属于 \(H^2\)，而\cref{b3:thm:28-bergman-coefficients}证明多项式在 \(A^2\) 中稠密，因此包含是稠密的。

为证明严格性，取具体的幂级数
\[
 F(z)=\sum_{n=0}^{\infty}\frac{z^n}{\sqrt{n+1}}.
\]
它在 \(\mathbb D\) 局部一致收敛，并满足
\[
 \|F\|_{A^2}^2=\sum_{n=0}^{\infty}\frac1{(n+1)^2}<\infty,
 \quad
 \sum_{n=0}^{\infty}\left|\frac1{\sqrt{n+1}}\right|^2
 =\sum_{n=0}^{\infty}\frac1{n+1}=\infty.
\]
所以 \(F\in A^2\)，但 \(F\notin H^2\)。它的 Taylor 多项式全部在 \(H^2\)，却在面积范数中趋于 \(F\)，同时证明 \(H^2\) 在 \(A^2\) 中不闭。
\end{proof}

这个反例还能直接排除 \(L^2\) 径向边界极限。它在半径 \(r\) 上的圆周范数满足
\[
 \int_0^{2\pi}|F(re^{\mathrm{i}\theta})|^2\,\mathrm{d}m(\theta)
 =\sum_{n=0}^{\infty}\frac{r^{2n}}{n+1}\longrightarrow\infty
 \quad(r\uparrow1).
\]
若这些径向函数在 \(L^2(\mathrm{d}m)\) 中收敛，其范数必有界，产生矛盾。另一方面，将上式按 \(2r\,\mathrm{d}r\) 积分，得到的正是收敛级数 \(\sum(n+1)^{-2}\)。因此面积平方可积与边界平方可积的差别，在同一个例子中表现为“平均值可以发散，但其径向积分仍有限”。

\ProseParagraphBreak

两个空间在共有函数上的范数也不等价：\(\|z^n\|_{H^2}=1\)，而 \(\|z^n\|_{A^2}=1/\sqrt{n+1}\to0\)。后续使用哪个核，必须由正在控制的范数决定。圆周边界问题使用 Hardy 核，面积正交投影与面积极值问题使用 Bergman 核；相同的解析函数在这两种几何结构下，可以具有不同的大小和不同的极值代表。


\begin{chaptersummary}
圆盘上的全纯性与积分可积性相结合，使点值成为连续线性泛函。对 \(H^2\)，圆周正交性给出系数平方和，再由同一组系数在圆周 \(L^2\) 中构造边界函数；Cauchy 表示和核 \((1-z\overline w)^{-1}\) 随之得到。对 \(A^2\)，径向积分给系数增加权重 \((n+1)^{-1}\)，标准正交基变成 \(\sqrt{n+1}\,z^n\)，再生核则为 \((1-z\overline w)^{-2}\)。核在对角上的值决定点值估计的最优常数，正交投影给出唯一的面积最佳逼近。

零点部分额外使用 Jensen 公式与边界范数。先取得 Blaschke 条件，再证明乘积在几乎处处的 \(L^2\) 边界上模长为一，才能完成等范数除零；无零点的余因子仍可能具有显著的边界奇性。范数之间也有方向性：\(H^2\) 连续嵌入 \(A^2\)，反向包含不成立；面积可积本身不提供平方可积的圆周边界。后续多变量 Bergman 空间将保留局部点值控制、完备性和核表示的思路，而不照搬一元零点因子分解。
\end{chaptersummary}

\BookExercises

% 题1—8均为按本章已证工具自拟或重新组织的训练；不是外部教材原题。
% 层次：边界反例；导数与插值；内函数子空间；紧性；算子计算；投影端点反例；自同构换元。

\begin{exercise}\label{b3:ex:28-1}
对整数 \(N\ge0\)，在单位圆盘上令
\[
 q_N(z)=\frac1{N+1}\sum_{n=0}^N z^n.
\]
计算 \(q_N\) 的 \(H^2\) 范数与归一化面积 \(A^2\) 范数。证明它在两种范数中均趋于零，但每个 \(q_N(1)=1\)。由此说明，在多项式上定义的边界求值 \(p\mapsto p(1)\) 不能连续延拓到这两个空间。再求 \(q_N(\zeta)\) 对每个 \(\zeta\in\mathbb T\setminus\{1\}\) 的极限，并解释为何不与 \(L^2\) 边界唯一性矛盾。
\end{exercise}
\begin{solution}
由两种系数范数公式，
\[
 \|q_N\|_{H^2}^2=\frac1{N+1},\qquad
 \|q_N\|_{A^2}^2=\frac1{(N+1)^2}\sum_{n=0}^N\frac1{n+1}
 \le\frac1{N+1}.
\]
所以两种范数都趋零，而代入一得到 \(q_N(1)=1\)。若边界求值存在连续线性延拓 \(E\)，则 \(|E(q_N)|\le C\|q_N\|\to0\)，与该恒等式矛盾。这里多项式在闭盘有通常点值，故反例甚至不涉及如何选择一般边界函数的代表。

当 \(|\zeta|=1\)、\(\zeta\ne1\) 时，有限几何和给出
\[
 |q_N(\zeta)|=\left|\frac{1-\zeta^{N+1}}{(N+1)(1-\zeta)}\right|
 \le\frac2{(N+1)|1-\zeta|}\longrightarrow0.
\]
因此边界逐点极限在一处为一，其余处为零。这个极限与零函数仅相差一个零测集，代表同一个 \(L^2(\mathbb T,m)\) 元素。边界的 \(L^2\) 唯一性本来只在几乎处处意义下成立，并没有指定每一点的值。
\end{solution}

\begin{exercise}\label{b3:ex:28-2}
设 \(w\in\mathbb D\)。在 \(H^2(\mathbb D)\) 上求导数泛函 \(f\mapsto f'(w)\) 的最优范数，证明
\[
 |f'(w)|\le\frac{\sqrt{1+|w|^2}}{(1-|w|^2)^{3/2}}\|f\|_{H^2}.
\]
求出所有取等函数，并在 \(w=0\) 处核对答案。
\end{exercise}
\begin{solution}
写 \(f(z)=\sum_{n\ge0}a_nz^n\)。对固定 \(w\)，序列 \((n w^{n-1})_{n\ge1}\) 平方可和，因而逐项导数值是绝对收敛的配对。令
\[
 d_w(z)=\sum_{n\ge1}n\overline w^{\,n-1}z^n
       =\frac{z}{(1-z\overline w)^2}.
\]
幂级数的逐项求导在紧子集上一致有效，故 \(f'(w)=\langle f,d_w\rangle_{H^2}\)。若 \(t=|w|^2<1\)，对几何级数两次求导可得
\[
 \|d_w\|_{H^2}^2=\sum_{n\ge1}n^2t^{n-1}
 =\frac{\mathrm d}{\mathrm dt}\frac{t}{(1-t)^2}
 =\frac{1+t}{(1-t)^3}.
\]
由 Cauchy--Schwarz 得所述上界。\(d_w\ne0\)，取 \(f=d_w\) 达到等号，所以常数最优；等号条件给出且仅给出 \(f=c d_w\)，其中 \(c\in\mathbb C\)，包括零函数。\(w=0\) 时 \(d_0(z)=z\)，结论成为 \(|a_1|\le(\sum|a_n|^2)^{1/2}\)，取等时仅一次项可能非零。
\end{solution}

\begin{exercise}\label{b3:ex:28-3}
设 \(a\in\mathbb D\setminus\{0\}\)、\(u,v\in\mathbb C\)。在满足 \(f(0)=u\)、\(f(a)=v\) 的全部 \(H^2(\mathbb D)\) 函数中，求范数最小者，证明其唯一，并计算最小范数平方。说明当 \(a\to0\) 而 \(u\ne v\) 固定时，答案为何发散。
\end{exercise}
\begin{solution}
先把常数项分离：\(f=u+h\)，其中 \(h(0)=0\)，且由系数正交性，
\[
 \|f\|_{H^2}^2=|u|^2+\|h\|_{H^2}^2.
\]
在无常数项的子空间中，评价 \(h(a)\) 由
\[
 v_a(z)=k_a(z)-1=\frac{z\overline a}{1-z\overline a},
 \qquad \|v_a\|^2=v_a(a)=\frac{|a|^2}{1-|a|^2}
\]
表示。令 \(d=v-u\)。Cauchy--Schwarz 给出
\(\|h\|^2\ge |d|^2/\|v_a\|^2\)。这个下界由
\[
 h_0=\frac{d}{\|v_a\|^2}v_a
\]
达到。任何其他满足条件的 \(h\) 都可写作 \(h_0+q\)，其中 \(q(0)=q(a)=0\)。于是 \(\langle q,v_a\rangle=q(a)=0\)，从而 \(\|h\|^2=\|h_0\|^2+\|q\|^2\)，证明唯一性。具体答案为
\[
 f_0(z)=u+\frac{(v-u)(1-|a|^2)}{a}\,
              \frac{z}{1-z\overline a},\qquad
 \|f_0\|^2=|u|^2+\frac{1-|a|^2}{|a|^2}|v-u|^2.
\]
当两个节点合并而函数值的差保持非零时，第二项趋于无穷。此时 \((f(a)-f(0))/a\) 被迫很大，与前题的局部导数控制相符。
\end{solution}

\begin{exercise}\label{b3:ex:28-4}
设 \(I\) 是 \(\mathbb D\) 上的内函数，即 \(I\) 全纯、\(|I|\le1\)，且其 \(L^2\) 边界函数满足 \(|I^*|=1\) 几乎处处。令
\[
 M=I H^2(\mathbb D),\qquad
 \mathcal K_I=\{f\in H^2:\langle f,Ih\rangle=0\text{ 对每个 }h\in H^2\}.
\]
证明 \(M\) 闭，且 \(\mathcal K_I\) 的再生核为
\[
 k_w^I(z)=\frac{1-\overline{I(w)}\cdot I(z)}{1-z\overline w}.
\]
计算它的范数平方。对 \(I(z)=z^m\)、整数 \(m\ge1\)，确定 \(\mathcal K_I\) 及其核；也处理 \(I\) 为模一常数的情形。
\end{exercise}
\begin{solution}
\cref{b3:prop:28-hardy-multiplier}说明乘以 \(I\) 是等距映射。若 \(Ih_j\) 在 \(H^2\) 中收敛，则 \(h_j\) 是 Cauchy 列，完备性给出 \(h_j\to h\)，故原极限为 \(Ih\)。这证明 \(M\) 闭。

对每个 \(w\in\mathbb D\)，把通常的 Hardy 核分成
\[
 k_w=\overline{I(w)}I k_w+k_w^I.
\]
第一项属于 \(M\)。对任意 \(h\in H^2\)，等距乘法也保持内积（或直接用边界表示），所以
\[
 \begin{aligned}
 \langle Ih,k_w^I\rangle
 &=I(w)\cdot h(w)-I(w)\langle Ih,I k_w\rangle\\
 &=I(w)\cdot h(w)-I(w)\langle h,k_w\rangle=0.
 \end{aligned}
\]
因而 \(k_w^I\in\mathcal K_I\)。若 \(f\in\mathcal K_I\)，它与第一项正交，故 \(f(w)=\langle f,k_w\rangle=\langle f,k_w^I\rangle\)。取 \(f=k_w^I\)，得到
\[
 \|k_w^I\|^2=k_w^I(w)=\frac{1-|I(w)|^2}{1-|w|^2}.
\]
当 \(I=z^m\) 时，\(M\) 恰由前 \(m\) 个 Taylor 系数为零的函数组成。因此
\[
 \mathcal K_I=\operatorname{span}\{1,z,\ldots,z^{m-1}\},
 \qquad k_w^I(z)=\sum_{j=0}^{m-1}z^j\overline w^{\,j}.
\]
这个有限几何和与核的通式相同。模一常数乘法满射，故此时 \(\mathcal K_I=\{0\}\)，核恒为零。这里仅研究给定内函数所生成的子空间，并未证明所有移位不变子空间都具有此形式。
\end{solution}

\begin{exercise}\label{b3:ex:28-5}
在圆周与面积均归一化为总质量一的约定下，证明包含映射 \(J:H^2(\mathbb D)\to A^2(\mathbb D)\) 是紧的：每个 \(H^2\) 范数有界序列都有在 \(A^2\) 范数中收敛的子列。再用单项式说明 \(H^2\) 的恒等映射没有这一性质，且从多项式的 \(A^2\) 范数控制其 \(H^2\) 范数不存在统一常数。
\end{exercise}
\begin{solution}
设 \(f_j=\sum_{n\ge0}a_{j,n}z^n\) 且 \(\sum_n|a_{j,n}|^2\le M^2\)。每个坐标有界，逐坐标取子列并作对角选择，得到一条子列，仍记作 \(f_j\)，使所有 \(a_{j,n}\) 都收敛。对每个固定 \(N\)，该子列的前 \(N+1\) 项在任何系数范数中都是 Cauchy 列；尾部则满足统一估计
\[
 \sum_{n>N}\frac{|a_{j,n}|^2}{n+1}\le\frac{M^2}{N+2}.
\]
对差 \(f_j-f_k\) 的尾部，使用 \(|a-b|^2\le2|a|^2+2|b|^2\)，得上界 \(4M^2/(N+2)\)。先令 \(N\) 足够大，再令 \(j,k\) 足够大，证明子列在 \(A^2\) 中是 Cauchy 列。该空间的完备性给出收敛。

相比之下，\(z^j\) 的 \(H^2\) 范数为一，任意不同两项的距离为 \(\sqrt2\)，故没有 \(H^2\) 收敛子列。最后
\[
 \|z^n\|_{H^2}=1,\qquad \|z^n\|_{A^2}=\frac1{\sqrt{n+1}},
\]
排除了反向控制常数。这也解释了紧性的来源：面积积分使高阶系数获得越来越小的权重，允许统一截断。
\end{solution}

\begin{exercise}\label{b3:ex:28-6}
在 \(A^2(\mathbb D)\) 中令 \(S f(z)=z\cdot f(z)\)，标准正交基记为 \(e_n(z)=\sqrt{n+1}\,z^n\)。求满足 \(\langle Sf,g\rangle=\langle f,S^*g\rangle\) 的伴随算子 \(S^*\) 在这组基上的作用。证明 \(S\) 的范数为一，但对每个非零 \(f\) 都有 \(\|Sf\|<\|f\|\)。计算 \(S^*S-SS^*\) 的全部对角系数，证明它是正算子，并可在算子范数中用有限秩算子逼近。
\end{exercise}
\begin{solution}
由基的归一化直接得到
\[
 Se_n=\sqrt{\frac{n+1}{n+2}}e_{n+1},\qquad
 S^*e_0=0,\qquad
 S^*e_n=\sqrt{\frac n{n+1}}e_{n-1}\quad(n\ge1).
\]
这些公式先在有限线性组合上满足伴随等式，且所定义算子的系数乘子均不超过一；因此它们连续延拓到全空间，等式也由稠密性保留。对 \(f=\sum c_ne_n\ne0\)，
\[
 \|Sf\|^2=\sum_{n\ge0}\frac{n+1}{n+2}|c_n|^2
 <\sum_{n\ge0}|c_n|^2.
\]
严格性来自至少一个 \(c_n\ne0\)。但 \(\|Se_n\|\to1\)，所以算子范数仍为一。逐个基向量计算给出
\[
 (S^*S-SS^*)e_n=\frac1{(n+1)(n+2)}e_n\qquad(n\ge0).
\]
故其二次型是 \(\sum |c_n|^2/((n+1)(n+2))\ge0\)。只保留 \(0\le n\le N\) 的对角系数得到有限秩算子，余项范数恰为 \(1/((N+2)(N+3))\)，趋于零。这个计算展示了与 Hardy 等距移位不同的面积权重效应。
\end{solution}

\begin{optionalexercise}\label{b3:ex:28-7}
在单位圆盘上令
\[
 F(w)=\frac{1-\overline w}{1-w}.
\]
证明 \(|F|=1\)，并计算它的 Bergman 投影 \(P F\)。证明 \(P F\) 在圆盘内无界，从而 \(P:L^2(\mathbb D,dA)\to A^2(\mathbb D)\) 的有界性不能加强为“每个有界输入都有有界解析投影”。
\end{optionalexercise}
\begin{solution}
圆盘内 \(1-w\ne0\)，分子为分母的共轭，所以 \(|F(w)|=1\)。特别地 \(F\in L^2\)，投影有定义。其第 \(n\) 个 Taylor 系数为
\[
 b_n=(n+1)\int_{\mathbb D}F(w)\cdot\overline w^{\,n}\,\mathrm dA(w).
\]
对每个固定 \(0<r<1\)，把 \((1-re^{\mathrm i\theta})^{-1}\) 展开成一致收敛的几何级数，再取角平均，得到
\[
 \int_{\mathbb T}F(r\zeta)\cdot\overline{r\zeta}^{\,n}\,\mathrm dm(\zeta)
 =r^{2n}-r^{2n+2}.
\]
在每个半径先做这一步，没有在整盘上声称几何级数一致收敛。之后的面积换序合法，因为 \(|F(w)\cdot\overline w^{\,n}|\le1\)。于是
\[
 b_n=(n+1)\int_0^1(r^{2n}-r^{2n+2})2r\,\mathrm dr
     =\frac1{n+2}.
\]
故
\[
 (P F)(z)=\sum_{n\ge0}\frac{z^n}{n+2}
 =\frac{-\log(1-z)-z}{z^2}\quad(z\ne0),
\]
其中对数取在 \(z=0\) 为零的解析分支，零处可去且值为 \(1/2\)。实数 \(r\uparrow1\) 时，分子趋于正无穷，故投影无界。同时 \(\sum_{n\ge0}1/((n+1)(n+2)^2)<\infty\)，与 \(P F\in A^2\) 完全相容；面积平方可积并不等于有界。
\end{solution}

\begin{optionalexercise}\label{b3:ex:28-8}
设 \(a\in\mathbb D\)，令
\[
 \varphi_a(z)=\frac{a-z}{1-\overline a z},\qquad
 q_a(z)=\frac{\sqrt{1-|a|^2}}{1-\overline a z}.
\]
证明 \(U_af=q_a(f\circ\varphi_a)\) 是 \(H^2(\mathbb D)\) 的等距满射，且 \(U_a^2\) 为恒等映射。说明权因子为何是圆周 Jacobi 因子模长的平方根，并与正文 \(A^2\) 的导数权因子比较。不得预先假定复合算子有界。
\end{optionalexercise}
\begin{solution}
\(\varphi_a\) 是闭盘附近解析的圆盘自同构，在圆周上为光滑双射。若 \(\varphi_a(e^{\mathrm i\theta})=e^{\mathrm i\beta(\theta)}\)，可局部选取光滑辐角；对恒等式求导，并代入 \(|e^{\mathrm i\theta}|=1\)，得
\[
 \beta'(\theta)=\frac{1-|a|^2}{|1-\overline a e^{\mathrm i\theta}|^2}
 =|q_a(e^{\mathrm i\theta})|^2>0.
\]
例如这个等式也可由 \(e^{\mathrm i\theta}\varphi'_a(e^{\mathrm i\theta})/\varphi_a(e^{\mathrm i\theta})\) 直接化简取得。由于映射为圆周的保向双射，\(\beta\) 在一周增加 \(2\pi\)，故圆周换元给出，对每个解析多项式 \(p\)，
\[
 \int_{\mathbb T}|q_a(\zeta)|^2\cdot|p(\varphi_a(\zeta))|^2\,\mathrm dm(\zeta)
 =\int_{\mathbb T}|p(\eta)|^2\,\mathrm dm(\eta).
\]
函数 \(U_ap\) 在闭盘附近全纯，其边界平方范数即 \(H^2\) 范数，因此 \(\|U_ap\|_{H^2}=\|p\|_{H^2}\)。

给定 \(f\in H^2\)，取多项式 \(p_j\to f\) 于 \(H^2\)。上述等距性使 \(U_ap_j\) 在 \(H^2\) 中收敛到某个函数。\(H^2\) 范数收敛蕴含紧子集上一致收敛，且 \(\varphi_a\) 把紧集映到紧集，所以该极限逐点等于 \(q_a(f\circ\varphi_a)\)。这既证明公式对所有 \(f\) 有意义，也证明等距性，没有预先使用复合算子的有界性。

直接计算
\[
 \varphi_a\circ\varphi_a(z)=z,\qquad
 1-\overline a\varphi_a(z)=\frac{1-|a|^2}{1-\overline a z},
 \qquad q_a(z)\cdot q_a(\varphi_a(z))=1,
\]
故 \(U_a^2f=f\)，从而满射。圆周长度换元的因子是 \(|\varphi'_a|\)，平方范数所需的解析权具有模 \(|\varphi'_a|^{1/2}\)；面积换元的因子是 \(|\varphi'_a|^2\)，所以 \(A^2\) 中用完整导数作解析权。两者不同的指数来自积分所在维数，并非核公式中任意选择的常数。
\end{solution}

% !TeX root = ../../Book3.tex
\OptionalChapter{多复变入门}{b3:ch:29}
\BookChapterNumber{b3:ch:29}
% 规划范围：A27–A31：建立模型上的多复变入口、局部 barpartial 与基本复流形语言；不建立一般 Levi 问题、全局 barpartial、层上同调或完整复几何。

\begin{chapterabstract}
待撰写。
\end{chapterabstract}

\begin{learninggoals}
待撰写。
\end{learninggoals}

% 本章导语与先修说明待撰写。
待撰写。

% !TeX root = ../../Book3.tex
\section{多变量全纯函数与 Cauchy 公式}
\label{b3:sec:2901}
% 规划课题：A27。

待撰写。

\subsection{复坐标与全纯映射}
\label{b3:subsec:290101}

待撰写。

\subsection{多圆盘上的 Cauchy 公式}
\label{b3:subsec:290102}

待撰写。

\subsection{幂级数表示与局部收敛}
\label{b3:subsec:290103}

待撰写。

\subsection{恒等原理与一元结论的边界}
\label{b3:subsec:290104}

待撰写。


% !TeX root = ../../Book3.tex
\section{分别全纯与 Hartogs 现象}
\label{b3:sec:2902}
% 规划课题：A27。

待撰写。

\subsection{分别全纯与联合全纯}
\label{b3:subsec:290201}

待撰写。

\subsection{Hartogs 延拓的模型区域}
\label{b3:subsec:290202}

待撰写。

\subsection{Hartogs 延拓的一个标准版本}
\label{b3:subsec:290203}

待撰写。

\subsection{多变量奇点与一元直觉}
\label{b3:subsec:290204}

待撰写。


% !TeX root = ../../Book3.tex
\section{多重次调和函数}
\label{b3:sec:2903}
% 规划课题：A28。
% 范围说明：中文主称在正式撰写时对照术语规范核定；本节对应 plurisubharmonic，不引入一般全局等价定理。

待撰写。

\subsection{复直线限制与次平均性质}
\label{b3:subsec:290301}

待撰写。

\subsection{上半连续性与基本运算}
\label{b3:subsec:290302}

待撰写。

\subsection{全纯函数对数的模型}
\label{b3:subsec:290303}

待撰写。

\subsection{光滑情形与 Levi 形式}
\label{b3:subsec:290304}

待撰写。


% !TeX root = ../../Book3.tex
\section{全纯域与伪凸性}
\label{b3:sec:2904}
% 规划课题：A28。

待撰写。

\subsection{全纯延拓与全纯域}
\label{b3:subsec:290401}

待撰写。

\subsection{伪凸性与穷竭函数}
\label{b3:subsec:290402}

待撰写。

\subsection{边界 Levi 条件}
\label{b3:subsec:290403}

待撰写。

\subsection{模型区域与全局问题的界线}
\label{b3:subsec:290404}

待撰写。


% !TeX root = ../../Book3.tex
\section{型分解与局部 \texorpdfstring{\(\bar\partial\)}{dbar} 方程}
\label{b3:sec:2905}
% 规划课题：A29。

待撰写。

\subsection{复微分形式的型分解}
\label{b3:subsec:290501}

待撰写。

\subsection{\texorpdfstring{\(\partial\)}{∂} 与 \texorpdfstring{\(\bar\partial\)}{dbar} 算子}
\label{b3:subsec:290502}

待撰写。

\subsection{\texorpdfstring{\(\bar\partial\)}{dbar} 闭性与相容条件}
\label{b3:subsec:290503}

待撰写。

\subsection{缩小多圆盘上的局部解}
\label{b3:subsec:290504}

待撰写。


% !TeX root = ../../Book3.tex
\section{多变量 Bergman 空间}
\label{b3:sec:2906}
% 规划课题：A30。

待撰写。

\subsection{有界域上的面积范数}
\label{b3:subsec:290601}

待撰写。

\subsection{局部点值估计与完备性}
\label{b3:subsec:290602}

待撰写。

\subsection{再生核的存在性}
\label{b3:subsec:290603}

待撰写。

\subsection{多圆盘上的核与正交展开}
\label{b3:subsec:290604}

待撰写。


% !TeX root = ../../Book3.tex
\section{复流形的基本语言}
\label{b3:sec:2907}
% 规划课题：A31。

待撰写。

\subsection{全纯图册与全纯映射}
\label{b3:subsec:290701}

待撰写。

\subsection{复切空间与型分解}
\label{b3:subsec:290702}

待撰写。

\subsection{复环面的坐标模型}
\label{b3:subsec:290703}

待撰写。

\subsection{复射影空间与复维一的情形}
\label{b3:subsec:290704}

待撰写。



\begin{chaptersummary}
待撰写。
\end{chaptersummary}

\BookExercises

待撰写。


